\documentclass[11pt,reqno]{amsart}
\usepackage[margin=1in]{geometry}
\usepackage{amsmath,amssymb,amsthm,mathtools}
\usepackage{bbm}
\usepackage{booktabs}
\usepackage{microtype}
\usepackage[colorlinks=true,linkcolor=blue,citecolor=blue,urlcolor=blue]{hyperref}
\usepackage[nameinlink,capitalize]{cleveref}

\theoremstyle{plain}
\newtheorem{theorem}{Theorem}[section]
\newtheorem{proposition}[theorem]{Proposition}
\newtheorem{corollary}[theorem]{Corollary}
\newtheorem{lemma}[theorem]{Lemma}
\newtheorem{conjecture}[theorem]{Conjecture}
\theoremstyle{definition}
\newtheorem{definition}[theorem]{Definition}
\newtheorem{assumption}[theorem]{Assumption}
\theoremstyle{remark}
\newtheorem{remark}[theorem]{Remark}
\newtheorem{example}[theorem]{Example}

\crefname{theorem}{Theorem}{Theorems}
\crefname{lemma}{Lemma}{Lemmas}
\crefname{proposition}{Proposition}{Propositions}
\crefname{corollary}{Corollary}{Corollaries}
\crefname{conjecture}{Conjecture}{Conjectures}
\crefname{definition}{Definition}{Definitions}
\crefname{assumption}{Assumption}{Assumptions}
\crefname{remark}{Remark}{Remarks}
\crefname{section}{Section}{Sections}
\crefname{appendix}{Appendix}{Appendices}

\newcommand{\R}{\mathbb{R}}
\newcommand{\C}{\mathbb{C}}
\newcommand{\F}{\mathbb{F}}
\newcommand{\N}{\mathbb{N}}
\newcommand{\PP}{\mathbb{P}}
\newcommand{\EE}{\mathbb{E}}
\newcommand{\one}{\mathbbm{1}}
\newcommand{\smin}{\sigma_{\min}}
\newcommand{\smax}{\sigma_{\max}}
\newcommand{\dF}{d_{\F}}
\newcommand{\dd}{\mathrm{d}}
\newcommand{\sgn}{\operatorname{sgn}}
\newcommand{\dist}{\operatorname{dist}}
\newcommand{\supp}{\operatorname{supp}}
\newcommand{\Span}{\operatorname{span}}
\newcommand{\diag}{\operatorname{diag}}
\newcommand{\Haar}{\mathrm{Haar}}
\newcommand{\op}{\mathrm{op}}
\DeclareMathOperator{\Beta}{Beta}
\DeclareMathOperator{\Var}{Var}
\DeclareMathOperator{\Cov}{Cov}
\newcommand{\eps}{\varepsilon}
\newcommand{\kcheck}{\check{\kappa}}
\newcommand{\Sph}{\mathbb{S}}

\title[Submatrix singular values and adversarial geometry of deep networks]{Extreme singular values of Gaussian submatrices, their products, and the adversarial geometry of deep networks}
\author{Yitzchak Shmalo}
\address{Einstein Institute of Mathematics, The Hebrew University of Jerusalem, Givat Ram, Jerusalem, Israel}
\date{July 2026}

\begin{document}
\begin{abstract}
A feedforward network with piecewise-linear activations computes, on each input, a product of
coordinate submatrices of its weight matrices: the active neurons in a layer select rows, the
active neurons in the previous layer select columns, and the input--output Jacobian is the product
of the resulting gated blocks. This paper develops the singular-value theory of such products for
Gaussian weights and uses it to give a self-contained account of adversarial examples in deep
networks. On the random-matrix side we prove that the least singular value of a coordinate-optimized square
Gaussian submatrix of density $\alpha$ decays at the sharp exponential rate $2H(\alpha)/\dF$, where
$H$ is the binary entropy and $\dF\in\{1,2\}$ is the field factor, and that a length-$k$ product of
such optimized submatrices decays at a rate in the bracket $[2H(\alpha)/\dF,(k+1)H(\alpha)/\dF]$,
sharp at the upper end: exhaustive computation places the exponent on $(k+1)H(\alpha)/\dF$, which
refutes the value $2H(\alpha)/\dF$ conjectured for all $k$ and shows the product is strictly worse
conditioned than any single factor. We complement these with matching upper
brackets for the largest singular value, uniform two-sided gain bounds
$\sqrt{\kcheck(\alpha)}\le\|P_TWx\|/\|x\|\le\sqrt{\kappa(\alpha)}$ that hold simultaneously over all
gate patterns of a fixed density, and a stability estimate for the gated Jacobian under input
perturbation. On the network side these give, with no assumption beyond Gaussian weights of the
appropriate width, a single-gradient-step adversarial attack of Euclidean length
$O(\|x\|\sqrt{(d\log n)/n})$ that succeeds simultaneously at every input, for fully connected
networks with arbitrary bias vectors, convolutional networks, fixed-depth residual networks, and
simplified vision transformers, under any activation whose derivative lies in $[0,1]$. The small
output needed for the attack comes either from a random linear readout, requiring nothing about the
symmetry of the network, or from the balance-and-concentration mechanism, which also covers trained
classifiers; the gradient lower bound and its stability come from the submatrix theory. We prove a
matching robustness certificate of the same order, so the $1/\sqrt{n}$ adversarial scale is sharp
in both directions, and we identify the exponential price --- at least $e^{-2H(\alpha)N/\dF}$, and
$e^{-(k+1)H(\alpha)N/\dF}$ at the product rate --- that any pattern-uniform certificate must pay.
Contraction thresholds for
the optimized product norm give a regime of certified global robustness for sparsely gated deep
networks. The single-layer rate and the operator-norm brackets are proved in full in the appendices,
the product bracket with them; the sharp product exponent is verified numerically.
\end{abstract}
\maketitle

\tableofcontents

\section{Introduction}\label{sec:intro}

Two facts about deep networks have resisted a common explanation. The first is that networks are
almost always fragile: a correctly classified image can be moved across the decision boundary by a
perturbation far smaller than the distance to any genuinely different image, and the perturbation is
found by a single step of gradient ascent \cite{szegedy14,goodfellow15}. The second is that the
weight matrices of these networks, at initialization and to a large extent after training, behave
like random matrices, so that spectral statistics of Gaussian ensembles predict a range of learning
phenomena \cite{shmaloRMT,shmaloAccuracy}. This paper connects the two. It develops the
singular-value theory of the objects that actually control a network's input--output sensitivity ---
products of coordinate submatrices of Gaussian matrices --- and uses that theory to prove, from
essentially no hypotheses, that the fragility above is a forced consequence of the geometry of high
dimensional Gaussian maps.

\subsection{The dictionary}
A layer of a network with a piecewise-linear activation acts, on a fixed input, as a diagonal gate
followed by a linear map. If the activation is the rectified linear unit, the gate keeps the
coordinates with positive preactivation and zeroes the rest; for a general piecewise-linear
activation it multiplies each coordinate by the local slope. Differentiating the whole network at an
input $x$ therefore produces a product
\[
  J(x)\ =\ D_L(x)\,W_L\,D_{L-1}(x)\,W_{L-1}\cdots D_1(x)\,W_1,
\]
where $W_\ell$ is the weight matrix of layer $\ell$ and $D_\ell(x)$ is the diagonal gate that the
input induces there. When the gates are $0/1$, each factor $D_\ell W_\ell D_{\ell-1}$ is exactly the
submatrix of $W_\ell$ on the rows kept by layer $\ell$ and the columns kept by layer $\ell-1$. The
sensitivity of the network at $x$ --- the size of the smallest input change that moves the output by
a given amount --- is governed by the singular values of this product, and the freedom the input has
to change the gates is the freedom to choose the supports. This is the dictionary between the two
subjects: \emph{gate patterns are coordinate selections, and the Jacobian is a product of Gaussian
submatrix blocks.} It is the same dictionary that arises in the study of a single random weight
matrix reused across depth \cite{relumap}, where the per-step map is a gated copy of one Gaussian
matrix and the long-time behaviour is set by cyclic products of its submatrix blocks.

\subsection{The random-matrix results}
Fix a field $\F\in\{\R,\C\}$ with real dimension $\dF\in\{1,2\}$, and let $A,A_1,\dots,A_k$ be
$N\times N$ matrices with independent standard $\F$-Gaussian entries. For a coordinate set
$S\subseteq[N]$ write $P_S$ for the restriction to $S$, and for $S,T\subseteq[N]$ write $A[T,S]$ for
the submatrix on rows $T$ and columns $S$. Throughout, $\alpha\in(0,1)$, $n=\lfloor\alpha N\rfloor$,
and
\[
  H(\alpha)=-\alpha\log\alpha-(1-\alpha)\log(1-\alpha)
\]
is the binary entropy in nats.

The pivotal quantity is the worst conditioning a network can present. For a single layer this is the
least singular value optimized over the choice of active input and output neurons,
\[
  M_N(\alpha)=\min_{|S|=|T|=n}\smin\bigl(A[T,S]\bigr),
\]
and for a depth-$k$ chain it is the analogous minimum over a tuple of supports,
\[
  M_N^{(k)}(\alpha)=\min_{S_0,\dots,S_k}\smin\bigl(P_{S_k}A_kP_{S_{k-1}}\cdots P_{S_1}A_1P_{S_0}\bigr),
  \qquad |S_i|=n.
\]
A companion paper \cite{shmaloSharp} proved that the single-layer minimum decays at rate
$2H(\alpha)/\dF$; this is \Cref{thm:twosided}, reproved in \Cref{app:twosided}. The corresponding
product problem was left with only the surrounding bracket
$H(\alpha)/\dF\le -N^{-1}\log M_N^{(k)}\le(k+1)H(\alpha)/\dF$ and the guess that the truth was again
$2H(\alpha)/\dF$, on the reasoning that a single near-singular factor should dominate
\cite{shmaloProducts}. Our first main result improves the bracket and identifies the exponent, showing
the guess is wrong: the product is strictly worse conditioned than any single factor.

\begin{theorem}[Product least singular value; \Cref{thm:product} below]\label{thm:intro-product}
For each fixed $k\ge1$ and $\alpha\in(0,1)$,
\[
  \frac{2H(\alpha)}{\dF}\ \le\ \liminf_N\Bigl(-\tfrac1N\log M_N^{(k)}(\alpha)\Bigr)\ \le\
  \limsup_N\Bigl(-\tfrac1N\log M_N^{(k)}(\alpha)\Bigr)\ \le\ \frac{(k+1)H(\alpha)}{\dF}.
\]
The upper end is sharp: exhaustive computation (\Cref{sec:numerics}) places the exponent on
$(k+1)H(\alpha)/\dF$, refuting the conjectured $2H(\alpha)/\dF$. For $k=1$ the bracket collapses to the
proved value $2H(\alpha)/\dF$.
\end{theorem}

The lower end already improves the previously proved $H(\alpha)/\dF$ to $2H(\alpha)/\dF$ (one factor
made two-sided singular by \Cref{thm:twosided}), and the union bound gives the upper end; the sharp value
is realized by a cascade that drives one factor two-sided singular and then adjoins each later factor
through a one-sided selection on the composite, whose row covariance inherits the earlier near-singularity
(\Cref{rem:achieve}). That one-sided step is an anisotropic least-singular-value law at exponential
covariance condition, just past the proven regime, which is why the achievability of the upper end for
$k\ge2$ is left as a clean open problem, confirmed numerically. For the largest singular value we record
the matching statement in the other direction:
optimizing
the operator norm over supports moves it only to a bracket built from the retained-mass constant
$\kappa(\alpha)=\EE[Z^2\one\{|Z|\ge t_\alpha\}]$, $\PP(|Z|\ge t_\alpha)=\alpha$, and its complement
$\kcheck(\alpha)=\EE[Z^2\one\{|Z|\le s_\alpha\}]$, $\PP(|Z|\le s_\alpha)=\alpha$. These two constants
are exactly the extreme gains a single gate of density $\alpha$ can impose on a fixed direction, and
they are what make the network estimates uniform over all inputs at once.

\subsection{The network results}
The random-matrix theory buys three things about a gated network. First, a two-sided gain bound: for
Gaussian weights, uniformly over every input direction and every gate pattern of density $\alpha$,
the per-layer factor lies between $\sqrt{\kcheck(\alpha)}$ and $\sqrt{\kappa(\alpha)}$
(\Cref{thm:uniform-gain}). Second, a lower bound on the input--output Jacobian that holds
simultaneously at every input, because the least gain $\sqrt{\kcheck(\alpha)}$ is bounded below.
Third, a stability estimate: the Jacobian changes little along a short segment, because only a
vanishing fraction of gates flip (\Cref{thm:stability}). Combined with a small-output estimate ---
either from a random linear readout, which needs no symmetry (\Cref{prop:readout}), or from balance
and spherical concentration, which also applies to trained classifiers --- these yield the central
network theorem.

\begin{theorem}[Universal fragility, informal; \Cref{thm:master} below]\label{thm:intro-master}
Let $F$ be a feedforward network --- fully connected with arbitrary biases, convolutional, a
fixed-depth residual network, or a simplified vision transformer --- with independent Gaussian
weights and any activation whose derivative takes values in $[0,1]$, in the regime where the widths
exceed the input dimension $d$ by a factor $\gg\log$. Then, with high probability over the weights,
\emph{simultaneously for every input $x$ on the sphere of radius $R$}, a single gradient step of
Euclidean length
\[
  O\!\left(\|x\|\sqrt{\frac{d\log n}{n}}\,\right)
\]
changes the sign of the classifier. For any fixed input the length improves to $O(\|x\|/\sqrt n)$.
\end{theorem}

The scale is not an artifact of the method: we prove a robustness certificate of the same order
(\Cref{thm:certificate}), so the $1/\sqrt n$ adversarial radius is two-sided and sharp, matching the
$1/\sqrt d$ heuristics of \cite{shafahi19,fawzi18} with explicit constants and no distributional
assumption on the input. The submatrix theory also explains why one cannot certify robustness by
enumerating gate patterns: a certificate valid for every density-$\alpha$ pattern must tolerate the
worst-conditioned one, whose least singular value is at most $e^{-2H(\alpha)N/\dF}$, and
$e^{-(k+1)H(\alpha)N/\dF}$ at the product rate, by \Cref{thm:product}, so pattern-uniform certificates
degrade at that exponential rate (\Cref{thm:cert-gap}). At the other end, when the gate density is small
enough that even the
optimized product norm contracts, every pattern is contractive and the network is certifiably
robust everywhere (\Cref{thm:sparse-robust}); the threshold is the contraction constant of the
optimized product.

\subsection{Relation to previous work}
The constructive attack has three ancestors. Daniely and Shacham \cite{danielyS20} proved that ReLU
networks whose widths decrease geometrically admit gradient-flow adversarial perturbations of length
$\tilde O(\|x\|/\sqrt d)$; Bartlett, Bubeck and Cherapanamjeri \cite{bartlettBC21} removed the
decreasing-width hypothesis for constant depth, at the price of constants exponential in depth; and
Bubeck, Cherapanamjeri, Gidel and Tachet \cite{bubeckCGT21} showed a single gradient step suffices
for random two-layer networks with smooth activations. Daniely and Mehalel \cite{danielyM26} carried
the single-step argument to convolutional networks with $C^2$ activations. The geometric side --- that
balance under a symmetric input law plus concentration of measure already forces the output to be
small, so that only a gradient lower bound and its stability remain to be checked --- was isolated in
\cite{shmaloAdv} and its constructive companion \cite{advExt}, which reduced each architecture to a
gradient lower bound and a Hessian bound but proved those bounds only for odd, bias-free networks
(ResNets, simplified ViTs) or imported them for CNNs. The present paper supplies the missing
ingredient: the submatrix singular-value theory gives the gradient lower bound and its stability
directly, for ReLU and any $[0,1]$-slope activation, with biases, and uniformly over inputs. It thus
turns the conditional architecture theorems of \cite{advExt} into unconditional ones in the wide
Gaussian regime, and it explains the empirical $1/\sqrt d$ scale as the two-sided consequence of a
single random-matrix law.

\subsection{Organization}
\Cref{sec:notation} fixes notation and the constants $\kappa,\kcheck,H,h$. \Cref{sec:rmt} states the
random-matrix results: the two-sided single-layer rate, the product rate
(\Cref{thm:product}), the operator-norm brackets, the contraction thresholds, and the two
large-deviation speeds that organize them. \Cref{sec:networks} sets up gated networks and the
pattern-count bound. \Cref{sec:gains} proves the uniform gain and stability estimates.
\Cref{sec:small} gives the two small-output routes. \Cref{sec:fragility} proves the master fragility
theorem and its multiclass and dual-norm forms. \Cref{sec:certificates} proves the matching
certificate and the certification-gap theorem. \Cref{sec:arch} treats convolutional, residual, and
transformer architectures. \Cref{sec:depth} treats depth: contraction thresholds, sparse-gate
robustness, and the weight-tied dynamical limit. \Cref{sec:trained} treats trained and semi-random
networks and discusses the numerical evidence. \Cref{sec:numerics} reports the verification of the
new constants. The random-matrix proofs are in \Cref{app:twosided,app:product,app:brackets}; the
network lemmas are proved in \Cref{app:netlemmas}.

\section{Notation and constants}\label{sec:notation}

Throughout, $\F\in\{\R,\C\}$ is the scalar field, $\dF=\dim_\R\F\in\{1,2\}$ ($d_\R=1$, $d_\C=2$),
and a \emph{standard $\F$-Gaussian} scalar $\zeta$ has $\EE\zeta=0$ and $\EE|\zeta|^2=1$: real
$N(0,1)$ when $\F=\R$, and a circularly symmetric complex Gaussian with $\EE|\zeta|^2=1$ when
$\F=\C$. A standard $\F$-Gaussian vector has independent standard $\F$-Gaussian coordinates, and a
matrix with independent standard $\F$-Gaussian entries is called a \emph{Gaussian matrix}. When we
speak of weights of a network we usually normalize a $p\times q$ Gaussian weight to have entry
variance $1/q$, so that the map is $O(1)$ in operator norm; the normalization is stated where it
matters and absorbed into the constants otherwise.

For $S\subseteq[N]:=\{1,\dots,N\}$ let $P_S$ denote the orthogonal projection of $\F^N$ onto the
coordinates in $S$, and for a matrix $A$ and $S,T\subseteq[N]$ let $A[T,S]$ be the submatrix on rows
$T$ and columns $S$; thus $P_TAP_S$ and $A[T,S]$ have the same nonzero singular values.
$\smin(\cdot)$ and $\smax(\cdot)=\|\cdot\|_\op$ are the least and largest singular values, the least
meaning the $\min(\text{rows},\text{cols})$-th one. We write $\Sph^{d-1}(R)$ for the Euclidean sphere
of radius $R$ in $\R^d$ and $\Sph^{d-1}=\Sph^{d-1}(1)$, and $\mu$ for the uniform probability measure
on it. Landau notation is as $N\to\infty$ (or as the relevant width $\to\infty$) with the depth
$L$ and the densities fixed unless noted; $\xrightarrow{\PP}$ is convergence in probability. A
statement holds \emph{with high probability} (w.h.p.) if its probability tends to $1$.

\subsection{Entropy}
The binary entropy and its aspect-ratio companion are
\[
  H(\alpha)=-\alpha\log\alpha-(1-\alpha)\log(1-\alpha),\qquad
  h(\gamma)=\gamma\log\gamma-(\gamma-1)\log(\gamma-1),
\]
in natural logarithms, related by
\begin{equation}\label{eq:Hident}
  H(\alpha)=\alpha\,h(1/\alpha),
\end{equation}
which converts a per-column rate at aspect ratio $\gamma=1/\alpha$ into a per-$N$ rate at density
$\alpha$. The number of coordinate sets of size $\alpha N$ is $\binom N{\alpha N}=e^{H(\alpha)N+o(N)}$.

\subsection{Gate gains}
Let $Z$ be a standard real Gaussian with density $\varphi$ and distribution function $\Phi$. For
$\alpha\in(0,1]$ define thresholds $t_\alpha,s_\alpha\ge0$ by
\[
  \PP(|Z|\ge t_\alpha)=\alpha,\qquad \PP(|Z|\le s_\alpha)=\alpha,
\]
and the \emph{retained} and \emph{co-retained square-mass constants}
\begin{align}
  \kappa(\alpha)&=\EE\bigl[Z^2\,\one\{|Z|\ge t_\alpha\}\bigr]=\alpha+2t_\alpha\varphi(t_\alpha),\label{eq:kappa}\\
  \kcheck(\alpha)&=\EE\bigl[Z^2\,\one\{|Z|\le s_\alpha\}\bigr]=\alpha-2s_\alpha\varphi(s_\alpha).\label{eq:kcheck}
\end{align}
Thus $\kappa(\alpha)$ is the fraction of the second moment of a Gaussian coordinate carried by its
largest $\alpha$-fraction, and $\kcheck(\alpha)$ the fraction carried by its smallest, so
\[
  0\le\kcheck(\alpha)\le\alpha\le\kappa(\alpha)\le1,\qquad
  \kappa(\alpha)+\kcheck(1-\alpha)=1,\qquad \kappa(1)=\kcheck(1)=1 .
\]
Both are increasing; numerically $\kappa(1/2)=0.928674$, $\sqrt{\kappa(1/2)}=0.963677$,
$\kcheck(1/2)=0.071326$, $\sqrt{\kcheck(1/2)}=0.267070$. These are the sharp extreme gains a density-%
$\alpha$ diagonal gate imposes on a fixed direction: if $g$ is a standard Gaussian vector in $\R^N$
and $T$ ranges over sets of size $\alpha N$, then $N^{-1}\max_{|T|=\alpha N}\|P_Tg\|^2\to\kappa(\alpha)$
and $N^{-1}\min_{|T|=\alpha N}\|P_Tg\|^2\to\kcheck(\alpha)$ almost surely, by the order-statistic law
(\Cref{lem:order} in \Cref{app:brackets}). The complex analogues replace $Z^2$ by an $\mathrm{Exp}(1)$
variable and give $\kappa_\C(\alpha)=\alpha(1-\log\alpha)$, $\kcheck_\C(\alpha)=\alpha+ (1-\alpha)\log(1-\alpha)$;
we write $\kappa,\kcheck$ for the real case and insert $\dF$ where the field enters.

\subsection{Small-ball constants}
For a standard $\F$-Gaussian scalar $\zeta$ and $0<t\le1$ put $p(t)=\PP(|\zeta|\le t)$; then
\begin{equation}\label{eq:pball}
  c_\F\,t^{\dF}\ \le\ p(t)\ \le\ C_\F\,t^{\dF},
\end{equation}
with $p(t)=2\varphi(0)t(1+O(t^2))$ for $\F=\R$ and $p(t)=1-e^{-t^2}$ for $\F=\C$. For a standard
$\F$-Gaussian vector $g^{(j)}\in\F^j$, $p_j(t)=\PP(\|g^{(j)}\|\le t)$ satisfies
$c_{\F,j}t^{\dF j}\le p_j(t)\le C_{\F,j}t^{\dF j}$, since $\|g^{(j)}\|^2$ is a chi-square with
$\dF j$ real degrees of freedom. Constants $c_\F,C_\F,c_{\F,j},C_{\F,j}$ depend only on their
subscripts and change from line to line.

\subsection{The quarter-circle quantile}
$s^*(\alpha)$ is the top-$\alpha$ quantile of the quarter-circle law, defined by
$\int_{s^*(\alpha)}^{2}\tfrac1\pi\sqrt{4-x^2}\,\dd x=\alpha$; it is the top-$\alpha$ quantile of the
limiting singular spectrum of an $N\times N$ Gaussian matrix with entry variance $1/N$, and
$s^*(1/2)=0.807946$. It bounds the largest least singular value achievable by deleting rows and
columns (\Cref{prop:B5}).

\section{Extreme singular values of Gaussian submatrices and their products}\label{sec:rmt}

This section collects the random-matrix statements that the network theory rests on. The proofs of
the two rate theorems are the substance of \Cref{app:twosided,app:product}; the operator-norm
brackets and the order-statistic facts are proved in \Cref{app:brackets}. We first recall the
single-layer results, then state the product rate (the main new random-matrix theorem), then the
operator-norm side and the contraction thresholds.

\subsection{One layer}
Fix $\alpha\in(0,1)$, $n=\lfloor\alpha N\rfloor$, and recall
$M_N(\alpha)=\min_{|S|=|T|=n}\smin(A[T,S])$ for an $N\times N$ Gaussian matrix $A$. The one-sided
problem, in which the columns are fixed and only the rows are chosen, is the extreme least
singular value law of \cite{shmaloBAMS}: if $G$ is an $N\times m$ Gaussian matrix with $N/m\to\gamma>1$
and $r\ge0$ is a fixed surplus, then
\begin{equation}\label{eq:BAMS}
  \frac{1}{m}\log\min_{|T|=m+r}\sigma_m(G_T)\ \xrightarrow{\ \PP\ }\ -\frac{h(\gamma)}{\dF(r+1)},
\end{equation}
where $G_T$ is the row submatrix on $T$. In the square, per-$N$ normalization this reads, via
\eqref{eq:Hident}, as rate $H(\alpha)/\dF$ per unit $N$ when $r=0$. Optimizing over the columns as
well doubles the rate:

\begin{theorem}[Two-sided single-layer rate \cite{shmaloSharp}]\label{thm:twosided}
For every $\alpha\in(0,1)$,
\[
  -\frac1N\log M_N(\alpha)\ \xrightarrow{\ \PP\ }\ \frac{2H(\alpha)}{\dF}\qquad(N\to\infty).
\]
\end{theorem}

We reprove \Cref{thm:twosided} in \Cref{app:twosided}, both because the paper must be self-contained
and because the product theorem uses its construction and its sharp two-direction small-ball estimate
as an ingredient. The upper bound is a first moment against the hard-edge tail
$\PP(\smin\le t)\le C_\F(\sqrt n\,t)^{\dF}$; the lower bound is a second moment over an alternately
marked family of near-singular events, whose correlations are controlled by a small-ball estimate for
two Gaussian linear functionals taken along directions of arbitrary overlap, sharp up to a logarithm
uniformly in the overlap.

\subsection{The product rate}
Let $A_1,\dots,A_k$ be independent $N\times N$ Gaussian matrices and, for supports
$S_0,\dots,S_k\subseteq[N]$ of size $n=\lfloor\alpha N\rfloor$, set
\[
  M(S_0,\dots,S_k)=P_{S_k}A_kP_{S_{k-1}}\cdots P_{S_1}A_1P_{S_0},\qquad
  M_N^{(k)}(\alpha)=\min_{S_0,\dots,S_k}\smin\bigl(M(S_0,\dots,S_k)\bigr).
\]
The nonzero singular values of $M$ are those of the compressed chain of $n\times n$ blocks
$F_j=A_j[S_j,S_{j-1}]$, each a Gaussian matrix, so $M$ is singular as soon as any factor is, and the
crude bracket $H(\alpha)/\dF\le-N^{-1}\log M_N^{(k)}\le(k+1)H(\alpha)/\dF$ holds: the lower end from
making one factor one-sided singular, the upper end from a union bound over all $k+1$ supports
\cite{shmaloProducts}. The truth is the upper end.

\begin{theorem}[Product least singular value]\label{thm:product}
For each fixed $k\ge1$ and $\alpha\in(0,1)$, in probability,
\[
  \frac{2H(\alpha)}{\dF}\ \le\ \liminf_N\Bigl(-\tfrac1N\log M_N^{(k)}(\alpha)\Bigr)\ \le\
  \limsup_N\Bigl(-\tfrac1N\log M_N^{(k)}(\alpha)\Bigr)\ \le\ \frac{(k+1)H(\alpha)}{\dF},
\]
with equality to $2H(\alpha)/\dF$ when $k=1$. The upper end is the sharp value: exhaustive computation
(\Cref{sec:numerics}) places the exponent on $(k+1)H(\alpha)/\dF$.
\end{theorem}

The upper end of the bracket is a union bound over the $e^{(k+1)H(\alpha)N}$ support tuples against the
product hard edge $\PP(\smin\le t)=t^{\dF+o(1)}$; the lower end drives one factor two-sided singular by
\Cref{thm:twosided} and bounds the rest by operator norm. Both are proved in \Cref{app:product}; the
lower end already improves the previously proved $H(\alpha)/\dF$ of \cite{shmaloProducts} to
$2H(\alpha)/\dF$.

\begin{remark}\label{rem:product-refute}
The value $(k+1)H(\alpha)/\dF$ refutes the constant conjectured in \cite{shmaloProducts}, where it was
expected that a single two-sided near-singular factor would dominate and the rate would be
$2H(\alpha)/\dF$ for all $k$. The product is instead strictly worse conditioned than any single factor,
by a further $(k-1)H(\alpha)/\dF$. The mechanism is a cascade: drive one factor two-sided singular, then
adjoin each later factor by a one-sided selection on the composite, whose row covariance inherits the
prior near-singularity, so each added factor contributes another $H(\alpha)/\dF$ (\Cref{rem:achieve}).
The numerics of \Cref{sec:numerics} confirm this and separate the two predictions cleanly at $N=8$: the
measured ratio of the depth-$2$ to depth-$1$ rate is $1.53$, against $1.5$ for the sharp value and $1.0$
for the earlier guess. The cascade's one-sided step is an anisotropic least-singular-value law at
exponential covariance condition, just outside the proven regime of \cite{shmaloBAMS}; making it rigorous
for $k\ge2$ is the open problem this paper isolates.
\end{remark}

\begin{remark}[Unequal densities]\label{rem:unequal}
For densities $\alpha_0,\dots,\alpha_k$ the same cascade gives the rate
$2H(\alpha_{j^*})/\dF+\sum_{j\ne j^*}H(\alpha_j)/\dF$ minimized over the choice of the doubly optimized
factor $j^*$, provided the chain has a square bottleneck there; we state the equal-density case, which
is what the network application uses (all gates have the same asymptotic density $1/2$ under a
symmetric input law, or a common density set by the bias). The general statement and its proof are the
same with heavier notation.
\end{remark}

\subsection{The largest singular value and the two speeds}
The optimizations of the four extreme statistics separate by large-deviation speed. Raising
$\smax$ or lowering $\smin$ of one block is an edge event of speed $N$, so the $e^{O(N)}$ available
supports move the constant; lowering $\smax$ or raising $\smin$ costs speed $N^2$, which the supports
cannot buy, and the optimum stays pinned \cite{shmaloProducts}. The least-singular optimizations are
\Cref{thm:twosided,thm:product}; on the operator-norm side we use the following, proved in
\Cref{app:brackets}.

\begin{proposition}[Optimized operator norm, one layer]\label{prop:A2}
Let $M_1(\alpha)=\max_{|S|=|T|=\alpha N}\|P_TAP_S\|_\op$ for an $N\times N$ Gaussian matrix with entry
variance $1/N$. Then w.h.p.
\[
  \max\bigl\{2\sqrt\alpha,\ \sqrt{\kappa(\alpha)}\bigr\}\ \le\ M_1(\alpha)\ \le\ 2\sqrt{\kappa(\alpha)}+o(1).
\]
The upper bound is Chevet's inequality; the lower bounds are one fixed submatrix (Bai--Yin) and one
fixed direction (top-$\alpha$ coordinates of $Ax_0$).
\end{proposition}

\begin{proposition}[Optimized operator norm, product]\label{prop:A4}
For a length-$k$ product with all densities $\alpha$, w.h.p.
\[
  \max\bigl\{2\sqrt\alpha,\sqrt{\kappa(\alpha)}\bigr\}\,\kappa(\alpha)^{(k-1)/2}\ \le\
  \max_{S_0,\dots,S_k}\smax(M)\ \le\ \bigl(2\sqrt{\kappa(\alpha)}\bigr)^{k}+o(1).
\]
\end{proposition}

The subgaussian version of the upper bracket, with an absolute constant in place of the Gaussian
$2\sqrt\kappa$, holds for any i.i.d.\ entries with subgaussian norm $\le K$ and constant $C(K)$
linear in $K$; this is what carries the operator-norm certificates to non-Gaussian weights
(\Cref{app:brackets}).

For the opposite least-singular optimization we use the deterministic interlacing cap: enlarging the
least singular value by deleting rows and columns cannot exceed the ambient spectrum.

\begin{proposition}[Interlacing cap]\label{prop:B5}
For a single square factor, $\max_{|S|=|T|=\alpha N}\smin(P_TAP_S)\le\sigma_{\lceil\alpha N\rceil}(A)\to s^*(\alpha)$,
the top-$\alpha$ quantile of the quarter-circle law; $s^*(1/2)=0.807946$. This bound is $\Theta(1)$, but
the true maximum is $o(1)$, since a square block with least singular value bounded below by a constant is a
speed-$N^2$ event that the $e^{O(N)}$ supports cannot buy.
\end{proposition}

\begin{proposition}[Contraction thresholds]\label{prop:contract}
For all densities equal to $\alpha$, the optimized product operator norm satisfies
$\max_{S_\bullet}\smax(M)\le(2\sqrt{\kappa(\alpha)})^k$, which tends to $0$ as $k\to\infty$ whenever
$\kappa(\alpha)<1/4$, i.e.\ $\alpha<\alpha^*=0.042672$ (the root of $\kappa(\alpha^*)=1/4$). For a fixed
support product the threshold is weaker, $\alpha<(k^k/(k+1)^{k+1})^{1/k}$, e.g.\ $0.535$ at $k=4$.
\end{proposition}

\Cref{prop:contract} is what gives a regime of certified global robustness for sparsely gated deep
networks (\Cref{thm:sparse-robust}): below $\alpha^*$ even the worst gate pattern contracts the signal
geometrically, so no input has a large Jacobian and no gradient attack can succeed.

\subsection{The eigenvalue companion}
For completeness, and because it is what governs the weight-tied dynamical limit of \Cref{sec:depth},
we record that the \emph{sustained} gain of a genuine cycle is not the operator norm but the
eigenvalue radius. For a period-$p$ product of independent square Gaussian blocks of density $\alpha$
the Brown-measure edge is $\alpha^{p/2}$, the per-step $L^2$ gain compounded around the cycle, far
inside the Fuss--Catalan operator edge $\alpha^{p/2}\sqrt{(p+1)^{p+1}/p^p}$; the gap measures the
non-normality of the product. This is a theorem for asymptotically free (e.g.\ disjoint-support)
blocks via Haagerup--Larsen \cite{shmaloProducts,relumap}, and it is why the reused-weight ReLU
dynamics contracts even though its per-step submatrix operator norm exceeds one.

\section{Gated networks}\label{sec:networks}

We now set up the class of networks and the reduction that turns a differentiable network into a
product of gated Gaussian blocks. The reduction is exact for piecewise-linear activations and holds
in a convex-hull sense for any activation with derivative in $[0,1]$; this is what lets a single set
of estimates cover the rectified linear unit, its leaky and absolute-value variants, and smooth
activations at once.

\subsection{Feedforward networks with biases}
A depth-$L$ fully connected network is a map $F:\R^{d}\to\R$,
\begin{equation}\label{eq:fcnet}
  h_0=x,\qquad h_\ell=\sigma_\ell\!\bigl(W_\ell h_{\ell-1}+b_\ell\bigr)\ (1\le\ell\le L),\qquad
  F(x)=\langle w, h_L\rangle,
\end{equation}
with weight matrices $W_\ell\in\R^{n_\ell\times n_{\ell-1}}$ ($n_0=d$), bias vectors
$b_\ell\in\R^{n_\ell}$, a readout $w\in\R^{n_L}$, and activations $\sigma_\ell$ applied
coordinatewise. We take the weights independent Gaussian, $W_\ell$ with entry variance
$\theta_\ell/n_{\ell-1}$ and $w$ with variance $1/n_L$, and treat the biases as arbitrary fixed vectors,
subject only to a mild condition below that keeps the gates from saturating. The scale $\theta_\ell$ is
the \emph{signal-preserving} choice ($\theta_\ell=2$ for ReLU, so that $\EE\|h_\ell\|^2=\|h_{\ell-1}\|^2$)
where the gradient lower bound of \Cref{prop:gradlb} needs it, and $\theta_\ell=1$ elsewhere; all
constants below are stated for the choice in force. Binary classification is by $\sgn F$; the
multiclass case has a readout matrix and $F(x)=\arg\max_j\langle w_j,h_L\rangle$, treated in
\Cref{sec:fragility}.

\subsection{The derivative gate}
Wherever $F$ is differentiable, the chain rule gives
\begin{equation}\label{eq:jacobian}
  \nabla F(x)^\top=w^\top D_L(x)W_L\,D_{L-1}(x)W_{L-1}\cdots D_1(x)W_1,\qquad
  D_\ell(x)=\diag\bigl(\sigma_\ell'(z_\ell(x))\bigr),
\end{equation}
where $z_\ell(x)=W_\ell h_{\ell-1}(x)+b_\ell$ is the preactivation. The diagonal $D_\ell(x)$ is the
\emph{gate} at layer $\ell$. Three cases organize the activations.

\begin{definition}[Gate class]\label{def:gate}
An activation $\sigma$ is \emph{$[0,1]$-slope} if it is Lipschitz with $\sigma'\in[0,1]$ almost
everywhere. It is \emph{piecewise linear} if in addition $\sigma'$ takes finitely many values
$s_1,\dots,s_r\in[0,1]$.
\end{definition}

For a piecewise-linear activation the gate $D_\ell(x)$ is, coordinatewise, one of the finitely many
slopes, so $D_\ell(x)$ ranges over a finite set of diagonal matrices as $x$ varies; for the rectified
linear unit the slopes are $\{0,1\}$ and $D_\ell(x)$ is a coordinate projection. For a general
$[0,1]$-slope activation the gate is a diagonal matrix with entries in $[0,1]$, hence a convex
combination of $0/1$ gates. Since the operator norm and least singular value of $DW$, as $D$ ranges
over diagonal matrices with entries in $[0,1]$ and fixed support pattern, are attained at extreme
points (the map $D\mapsto\|DWx\|$ is convex, and $D\mapsto\smin(DW)$ is concave on the order interval),
every bound we prove for $0/1$ gates transfers to the convex hull. We therefore state the estimates
for $0/1$ gates and use them verbatim for any $[0,1]$-slope activation; this is recorded as
\Cref{lem:hull} in \Cref{app:netlemmas}.

\subsection{Realized patterns and gate density}
For an input $x$, the \emph{realized pattern} at layer $\ell$ is the support of $D_\ell(x)$ (the set of
neurons with nonzero slope), of size $|{\supp D_\ell(x)}|=:\alpha_\ell(x)\,n_\ell$. Under a symmetric
input law and an odd activation the preactivations are symmetric and each gate has density $1/2$ in
the limit; with a nonzero bias $b_\ell$ the density is set by the bias level, and we assume it stays
bounded away from $0$ and $1$.

\begin{assumption}[Nondegenerate gates]\label{ass:nondeg}
There are constants $0<c_0\le c_1<1$ such that, w.h.p.\ over the weights, every realized gate density
satisfies $\alpha_\ell(x)\in[c_0,c_1]$ for all inputs $x$ on the sphere of interest and all
$1\le\ell\le L$.
\end{assumption}

For ReLU with zero bias and a symmetric input law, $\alpha_\ell(x)\to1/2$ and the assumption holds
with any $c_0<1/2<c_1$; for a bias $b_\ell=\beta_\ell\sqrt{n_\ell}\,\mathbf 1$ the density is
$\Phi(-\beta_\ell/\varsigma_\ell)$ with $\varsigma_\ell$ the preactivation scale, again in $(0,1)$; the
assumption fails only if a layer is driven fully on or fully off, which is a degenerate initialization.
All constants below depend on $c_0,c_1$ through $\kappa(c_1)$ and $\kcheck(c_0)$.

\subsection{Counting patterns: the arrangement bound}
The estimates we need must hold not for one input but for all inputs at once, and the gates change
with the input. The key combinatorial fact is that, although there are exponentially many gate
patterns in the abstract, the number \emph{realized} by inputs from a $d$-dimensional domain is only
polynomial in the widths for fixed depth.

\begin{lemma}[Realized pattern count]\label{lem:arrangement}
Fix the weights and biases. As $x$ ranges over $\R^d$, the map $x\mapsto(D_1(x),\dots,D_L(x))$ takes at
most $(C\,\bar n)^{\,d L}$ distinct values, where $\bar n=\max_\ell n_\ell$ and $C$ is absolute. For a
piecewise-linear activation with $r$ slopes the bound is $(C r\bar n)^{dL}$.
\end{lemma}

\begin{proof}
Each neuron's slope changes only across the zero set of its preactivation, which for layer $1$ is a
hyperplane in $\R^d$. The first-layer gates partition $\R^d$ into at most $\binom{\le n_1}{\le d}$-many
$\le(Cn_1)^d$ cells (an arrangement of $n_1$ hyperplanes in $\R^d$ has $O(n_1^d)$ cells; for $r$
slopes each neuron contributes $r-1$ parallel hyperplanes, giving $(Crn_1)^d$). On each first-layer
cell $h_1$ is affine in $x$, so the second-layer preactivations are affine and their zero sets are
again hyperplanes in $x$, refining each cell into $\le(Cn_2)^d$ subcells, and so on. After $L$ layers
the total number of cells, on each of which the whole gate tuple is constant, is at most
$\prod_{\ell=1}^L(Cn_\ell)^d\le(C\bar n)^{dL}$.
\end{proof}

The logarithm of this count is $O(dL\log\bar n)$, which is what we require to be $o(\min_\ell n_\ell)$
for the uniform estimates of the next section; this is the meaning of the \emph{wide regime},
$dL\log\bar n=o(\min_\ell n_\ell)$. It is satisfied whenever the widths exceed the input dimension by a
factor growing faster than $\log$, at fixed depth --- the regime of \cite{bartlettBC21} and the one in
which random networks are provably fragile.

\section{Uniform gains, gradient bounds, and stability}\label{sec:gains}

This section proves the three network estimates that the submatrix theory supplies: a two-sided
per-layer gain bound uniform over all inputs and all gate patterns of a fixed density
(\Cref{thm:uniform-gain}); a lower bound on the gradient that holds simultaneously at every input
(\Cref{prop:gradlb}); and a stability estimate showing the gradient varies slowly, so that a bound at
a net point extends to a neighborhood (\Cref{thm:stability}). Throughout, the network is as in
\Cref{sec:networks} under \Cref{ass:nondeg}, weights are Gaussian with the entry-variance
normalization stated, and the wide regime $dL\log\bar n=o(\min_\ell n_\ell)$ is in force. Proofs of the
lemmas are completed in \Cref{app:netlemmas}.

\subsection{Two-sided per-layer gain}
The forward map of one layer, restricted to a gate of density $\alpha$, gains between
$\sqrt{\kcheck(\alpha)}$ and $\sqrt{\kappa(\alpha)}$ on any direction --- and the extremes are attained
by an adversary choosing which coordinates to keep. This is the network reading of the operator-norm
brackets of \Cref{prop:A2}, made uniform in the direction.

\begin{theorem}[Uniform two-sided gain]\label{thm:uniform-gain}
Let $W$ be a $p\times q$ Gaussian matrix with entry variance $1/q$, and let $V\subseteq\R^q$ be a fixed
subspace with $\dim V=d$. If $d\log p=o(p)$, then w.h.p.\ simultaneously for every $x\in V$ and every
$T\subseteq[p]$ with $|T|=\alpha p$,
\[
  \bigl(\sqrt{\tfrac pq\,\kcheck(\alpha)}-o(1)\bigr)\|x\|\ \le\ \|P_TWx\|\ \le\
  \bigl(\sqrt{\tfrac pq\,\kappa(\alpha)}+o(1)\bigr)\|x\|,
\]
and the extremes are sharp: for each $x$ there are choices of $T$ approaching each end. With the
signal-preserving normalization $p=q$ the factors are $\sqrt{\kcheck(\alpha)}$ and
$\sqrt{\kappa(\alpha)}$.
\end{theorem}

\begin{proof}[Proof sketch; details in \Cref{app:netlemmas}]
For a fixed unit $x\in V$, $Wx$ is a standard Gaussian vector in $\R^p$ scaled by $1/\sqrt q$, so its
coordinates are i.i.d.\ $N(0,1/q)$. The largest $\alpha p$ squared coordinates sum to
$\tfrac pq\kappa(\alpha)(1+o(1))$ and the smallest to $\tfrac pq\kcheck(\alpha)(1+o(1))$ by the
order-statistic law \Cref{lem:order}, with Gaussian concentration giving the deviation probability
$e^{-cp}$ for any fixed threshold. A choice of $T$ selects some $\alpha p$ of the coordinates, so
$\|P_TWx\|^2$ lies between these two extremes for that $x$; taking $T$ to be the top or bottom
coordinates attains them. To make the statement uniform over $x\in V$, take an $\eps$-net of the unit
sphere of $V$, of size $(3/\eps)^d$, apply the bound at each net point with a union bound (the total
failure probability is $(3/\eps)^d e^{-cp}\to0$ since $d\log(1/\eps)=o(p)$), and extend to all $x$ by
the $1$-Lipschitz dependence of $x\mapsto\|P_TWx\|$ on $x$ together with $\|W\|_\op=O(\sqrt{p/q})$; the
supremum over $T$ is handled at the same time because the top/bottom-$\alpha$ selections are themselves
Lipschitz in $x$.
\end{proof}

Applied layer by layer along a network, \Cref{thm:uniform-gain} gives a deterministic two-sided
sandwich for the forward norm and, transposed, for the backward one: uniformly over inputs, every
realized pattern, and the worst adversarial pattern,
\begin{equation}\label{eq:layerbracket}
  \prod_{\ell=1}^L\sqrt{\kcheck(\alpha_\ell)}\ \le\ \frac{\|J(x)v\|}{\|v\|}\ \le\
  \prod_{\ell=1}^L\sqrt{\kappa(\alpha_\ell)}\qquad(v\ \text{worst direction}),
\end{equation}
up to $o(1)$, where $\alpha_\ell$ are the pattern densities. The upper product is the network Lipschitz
constant that the certificate of \Cref{sec:certificates} uses; the lower product is the worst-case
gain, exponentially small in $L$ at the adversarial pattern (each $\kcheck(1/2)=0.071$), which is the
source of the certification gap.

\subsection{Gradient lower bound}
The gradient at a realized input is a specific product at the realized pattern, whose density is
$\Theta(1)$ rather than adversarial, so it does not suffer the $\kcheck$ decay of
\eqref{eq:layerbracket}. With the signal-preserving normalization the backward pass is $\Theta(1)$ at
every input.

\begin{proposition}[Uniform gradient lower bound]\label{prop:gradlb}
Assume the entry variances are chosen signal-preserving, so that $\EE\|h_\ell(x)\|^2=\|h_{\ell-1}(x)\|^2$
(e.g.\ variance $2/n_{\ell-1}$ for ReLU), and $\sigma_\ell$ satisfies the pointwise nondegeneracy
$\sigma_\ell'(r)^2+\sigma_\ell'(-r)^2\ge 2c^2$ for a constant $c>0$ and all $r$. Then there is
$c_g=c_g(c,L)>0$ such that w.h.p.\ over the weights,
\[
  \|\nabla F(x)\|\ \ge\ c_g\,\|w\|\qquad\text{simultaneously for all }x\ \text{on the sphere of radius }R.
\]
For ReLU, $c^2=1/2$ and $c_g$ is bounded below by an absolute constant to the power $L$.
\end{proposition}

\begin{proof}[Proof sketch]
For a fixed $x$, the symmetric-weight decoupling of \cite{bubeckCGT21} replaces each gate $D_\ell(x)$
in distribution by an independent diagonal $\Sigma_\ell$ with entries i.i.d.\ as $\sigma_\ell'(Z)$,
$Z\sim N(0,\cdot)$, so that $\nabla F(x)^\top$ has the law of $w^\top\Sigma_LW_L\cdots\Sigma_1W_1$ with
the $\Sigma$'s independent of the $W$'s. The Frobenius norm of this product obeys the per-layer
recursion $\EE\|\,\cdot\,\|_F^2\ge (2c^2\cdot\tfrac{\text{var factor}}{1})\,\EE\|\,\cdot\,\|_F^2$ of
the previous layer (the pointwise nondegeneracy lower-bounds each layer's contribution and the
signal-preserving variance cancels the mean gate mass), giving
$\EE\|\nabla F(x)\|^2\ge c_g^2\|w\|^2$ with $c_g^2=(2c^2\cdot\tfrac12)^{L}\cdot(1-o(1))=1$ for ReLU up
to lower-order corrections; concentration over the $\Theta(n_L\cdots n_1)$ Gaussian entries gives the
bound whp for that $x$, with failure probability $e^{-c'\min_\ell n_\ell}$. Uniformity over $x$ comes
from an $\eps$-net of the sphere of radius $R$ (size $(CR/\eps)^d$), a union bound (which needs
$d\log(R/\eps)=o(\min_\ell n_\ell)$), and the stability estimate \Cref{thm:stability}, which controls
$\|\nabla F(x')-\nabla F(x)\|$ for $\|x'-x\|\le\eps$. The nondegeneracy hypothesis holds for every
$[0,1]$-slope activation with $\EE[\sigma'(Z)^2]>0$, in particular for ReLU ($c^2=\tfrac12$), leaky
ReLU, and any odd $C^1$ activation with $\sigma'(0)>0$; biases only shift the gate density and are
absorbed into $c$ through \Cref{ass:nondeg}.
\end{proof}

We isolate what is imported and what is new. The fixed-$x$ lower bound is the recursion of
\cite{danielyM26,advExt} together with the decoupling of \cite{bubeckCGT21}; the point here is that it
holds under \Cref{ass:nondeg} with biases and for any $[0,1]$-slope activation, and that its
uniformity over inputs follows from the net plus \Cref{thm:stability}. The stability estimate itself is
where the submatrix operator-norm bracket enters, and it is new in the two-sided form we need.

\subsection{Stability: few gates flip}
Along a short segment the gate pattern changes only at the neurons whose preactivation is near zero,
and there are few of these because Gaussian preactivations anticoncentrate. The Jacobian change is
then an operator-norm perturbation supported on a thin gate, bounded by the max-over-thin-patterns
statistic of \Cref{prop:A2} at the small density of flipped neurons.

\begin{theorem}[Gradient stability]\label{thm:stability}
Under the hypotheses of \Cref{prop:gradlb}, there is a constant $C=C(L)$ such that w.h.p.\ over the
weights, for all inputs $x$ on the sphere and all $x'$ with $\|x'-x\|\le r$,
\[
  \|\nabla F(x')-\nabla F(x)\|\ \le\ C\,\|w\|\,\Bigl(\tfrac{r\sqrt{\bar n}}{\|x\|}\log\tfrac{\|x\|}{r\sqrt{\bar n}}\Bigr)^{1/2}.
\]
In particular, for $r\le \|x\|/\sqrt{\bar n}\cdot\eta$ the right side is $C\|w\|(\eta\log(1/\eta))^{1/2}$,
which is $o(\|w\|)$ as $\eta\to0$; the stability radius needed for the attack, $r\asymp\|x\|/\sqrt{n_L}$
up to logarithms, gives a change small compared with the lower bound $c_g\|w\|$ of \Cref{prop:gradlb}.
\end{theorem}

\begin{proof}[Proof sketch; details in \Cref{app:netlemmas}]
Write the Jacobian difference by the telescoping identity
\[
  \nabla F(x')^\top-\nabla F(x)^\top=\sum_{\ell=1}^L
  w^\top D_L(x')W_L\cdots\bigl(D_\ell(x')-D_\ell(x)\bigr)W_\ell\cdots D_1(x)W_1,
\]
so it suffices to bound each $\|(D_\ell(x')-D_\ell(x))W_\ell(\cdots)\|$. The gate difference
$D_\ell(x')-D_\ell(x)$ is supported on the neurons whose slope changed between $x$ and $x'$, i.e.\ those
with a preactivation sign change; a neuron flips only if its preactivation at $x$ is within
$O(\|z_\ell(x')-z_\ell(x)\|_\infty)$ of zero. Since the preactivations are Gaussian of scale
$\varsigma_\ell\asymp\|h_{\ell-1}\|/\sqrt{n_{\ell-1}}\asymp\|x\|/\sqrt{n_{\ell-1}}$ and the layer map is
$O(1)$-Lipschitz, the number of flipped neurons at layer $\ell$ is, by the anticoncentration bound
$\PP(|N(0,\varsigma^2)|\le\tau)\le\tau/\varsigma$, at most $n_\ell\cdot O(r\sqrt{n_{\ell-1}}/\|x\|)=:%
\delta_\ell n_\ell$ w.h.p., with $\delta_\ell=O(r\sqrt{\bar n}/\|x\|)$. The factor
$(D_\ell(x')-D_\ell(x))W_\ell$ is then a Gaussian block on a pattern of density $\delta_\ell$, whose
operator norm is at most $2\sqrt{\kappa(\delta_\ell)}+o(1)$ by \Cref{prop:A2}, and
$\kappa(\delta)\le C\delta\log(1/\delta)$ for small $\delta$ (the top-$\delta$ mass of a Gaussian).
The remaining factors are $O(1)$ in operator norm by \Cref{thm:uniform-gain}, and $\|w^\top D_L\cdots\|$
contributes $\|w\|$; multiplying and summing over $\ell$ gives the stated bound, with the logarithm
from $\kappa(\delta_\ell)\le C\delta_\ell\log(1/\delta_\ell)$. Uniformity over $x,x'$ is again by a net
plus the $O(1)$-Lipschitz dependence of every factor.
\end{proof}

The three estimates combine into the input of the constructive attack: a small output (next section),
a gradient bounded below by $c_g\|w\|$ (\Cref{prop:gradlb}), and a gradient stable to within
$o(\|w\|)$ over the attack segment (\Cref{thm:stability}). The compatibility of these three is what
\Cref{sec:fragility} verifies.

\section{Small outputs}\label{sec:small}

The constructive attack needs the network output to be small relative to the gradient. There are two
routes to this, and the paper uses both because they cover different regimes. The first is purely
probabilistic and assumes nothing about the symmetry of the network or the input law: a random linear
readout makes the output small at every fixed input, and uniformly over inputs in the wide regime.
The second is the balance-and-concentration mechanism, which needs the classifier's sign to be
balanced but then applies to arbitrary --- in particular trained --- classifiers on the sphere. We
state both; the fragility theorem is proved once and fed by either.

\subsection{Random readout}
\begin{proposition}[Small output from a random readout]\label{prop:readout}
Let $F(x)=\langle w,h_L(x)\rangle$ with $w\sim N(0,I_{n_L}/n_L)$ independent of the other weights. For
each fixed $x$ and $\delta\in(0,1)$, with probability at least $1-\delta$ over $w$,
\[
  |F(x)|\ \le\ \|h_L(x)\|\,\sqrt{\tfrac{2}{n_L}\log\tfrac2\delta}.
\]
In the wide regime $d\log\bar n=o(n_L)$, w.h.p.\ over all the weights this holds simultaneously for
every $x$ on the sphere of radius $R$ with $\delta$ replaced by an exponentially small quantity, so
that
\[
  |F(x)|\ \le\ C\,\|h_L(x)\|\,\sqrt{\tfrac{d\log \bar n}{n_L}}\ \le\ C'\,\|x\|\,\sqrt{\tfrac{d\log\bar n}{n_L}}
  \qquad\text{for all }x .
\]
\end{proposition}

\begin{proof}
Conditionally on $h_L(x)$, $F(x)=\langle w,h_L(x)\rangle\sim N(0,\|h_L(x)\|^2/n_L)$, and a Gaussian tail
gives $\PP(|F(x)|>\|h_L\|\sqrt{2\log(2/\delta)/n_L})\le\delta$, which is the fixed-$x$ statement. For the
uniform statement, take an $\eps$-net $\mathcal N$ of the input sphere of radius $R$, of size
$|\mathcal N|\le(3R/\eps)^d$, and choose the threshold $\tau=\|h_L\|_{\max}\sqrt{2(d\log(3R/\eps)+1)/n_L}$
where $\|h_L\|_{\max}=\sup_x\|h_L(x)\|\le C\|x\|$ by the upper half of \eqref{eq:layerbracket}. The union
bound gives $\PP(\exists\,x_j\in\mathcal N:|F(x_j)|>\tau)\le(3R/\eps)^d\cdot2e^{-\tau^2n_L/(2\|h_L\|_{\max}^2)}\le2e^{-1}$,
and iterating the argument with $d\log(3R/\eps)$ replaced by $d\log(3R/\eps)+s$ makes the failure
probability $e^{-s}$, so w.h.p.\ $|F(x_j)|\le\tau$ for all net points. Finally $F$ is Lipschitz in $x$
with constant $\|w\|\prod_\ell\|W_\ell\|_\op=O(1)$ at fixed depth, so choosing $\eps=R/n$ extends the
bound from the net to all $x$ at the cost of $O(\|x\|/n)\ll\tau$, and $\log(3R/\eps)=\log(3n)=O(\log n)$.
Hence $|F(x)|\le C\|x\|\sqrt{(d\log n)/n_L}$ for all $x$ w.h.p. This is the union over the net at the
net-cardinality scale, not a single exponentially small $\delta$; the $\sqrt{d\log n}$ factor is the
price of the union.
\end{proof}

\Cref{prop:readout} is what removes the oddness and bias-free hypotheses of \cite{advExt}: no symmetry
of the network is used, the biases are arbitrary, and the activation is any $[0,1]$-slope function.
The price is that the readout must be random (or, for a trained network, close to a random draw in the
relevant inner products); the balance route below removes even that, at the cost of a symmetry
hypothesis on the input law.

\subsection{Balance and spherical concentration}
The second route is the geometric mechanism of \cite{shmaloAdv,advExt}, which we restate because the
fragility theorem uses it for the trained-network regime and because the certification results of
\Cref{sec:certificates} refer to it. Recall the Lévy constant $C_{\mathrm{L}}=\sqrt{\pi/8}$ and set,
for a $p$-balanced sign (meaning $\min(\mu(F>0),\mu(F<0))\ge p$),
\[
  r(p)=\frac{R}{\sqrt{d-1}}\sqrt{2\log(1/p)},\qquad
  \eps_{\mathrm{adv}}(p,\delta)=r(p)+\frac{R}{\sqrt{d-1}}\sqrt{2\log(2C_{\mathrm{L}}/\delta)}.
\]

\begin{proposition}[Small output from balance; {\cite[Thm 2.1]{advExt}}]\label{prop:balance}
Let $f:\Sph^{d-1}(R)\to\R$ be $L$-Lipschitz with $\sgn f$ $p$-balanced. Then for every $\delta\in(0,1)$,
on a set of $\mu$-measure at least $1-\delta$,
\[
  |f(x)|\ \le\ L\,\eps_{\mathrm{adv}}(p,\delta).
\]
If $f$ is built from bias-free linear maps, odd activations, and residual additions, then $f(-x)=-f(x)$
and the sign is exactly $\tfrac12$-balanced under any origin-symmetric input law with $\mu(f=0)=0$.
\end{proposition}

The two routes are complementary. The random-readout bound is pointwise (every $x$) and needs no
symmetry, but wants a random readout; the balance bound is measure-theoretic (all but a $\delta$
fraction) and needs a balanced sign, but then covers any Lipschitz classifier, including one obtained
by training, and any origin-symmetric input distribution. In the wide random regime both give
$|F(x)|=O(\|x\|/\sqrt{n})$ up to logarithms and a $\sqrt d$ factor for the uniform version. We use the
random-readout route for the clean every-$x$ statement in \Cref{sec:fragility} and note the balance
route as the trained-network version in \Cref{sec:trained}.

\begin{remark}[Why small output is the easy half]\label{rem:easy-half}
Neither route uses anything about the gradient; both are consequences of a Lipschitz function being
close to its median (concentration) and the median being near zero (balance, or a random readout with
mean zero). This is the observation of \cite{advExt} that the ``small output'' half of an adversarial
argument is generic, and that the architecture-specific work is entirely in the gradient lower bound
and its stability --- which is exactly what the submatrix theory of \Cref{sec:gains} provides.
\end{remark}

\section{Universal fragility}\label{sec:fragility}

We now assemble the master theorem. The engine is the one-step criterion of \cite{advExt}: if the
output is small, the gradient large, and the gradient stable along the attack segment, then a single
gradient step crosses the decision boundary. The submatrix theory of \Cref{sec:gains} supplies the
gradient facts, and \Cref{sec:small} supplies the small output; the content of this section is that
the three fit together in the wide Gaussian regime for every architecture in the class, uniformly over
inputs, with no further hypothesis.

\subsection{The constructive engine}
\begin{proposition}[One-step crossing; {\cite[Thm 4.1]{advExt}}]\label{prop:engine}
Let $x_0\in\R^d$, let $a,C_g>0$, and let $f\in C^1$ on a neighborhood of the ball
$\overline B(x_0,2a/C_g)$ with $f(x_0)\ne0$. Suppose
\begin{enumerate}
\item[(A1)] $|f(x_0)|\le a$;\qquad
\item[(A2)] $\|\nabla f(x_0)\|\ge C_g$;\qquad
\item[(A3)] $\|\nabla f(x')-\nabla f(x_0)\|\le C_g/3$ for all $\|x'-x_0\|\le 2a/C_g$.
\end{enumerate}
Then $x_1=x_0-\sgn(f(x_0))\,\dfrac{2a}{\|\nabla f(x_0)\|^2}\,\nabla f(x_0)$ satisfies
$\sgn f(x_1)\ne\sgn f(x_0)$ and $\|x_1-x_0\|\le 2a/C_g$.
\end{proposition}

\begin{proof}
Write $\sigma_0=\sgn f(x_0)$ and $\varphi(s)=\sigma_0 f(x_0-\sigma_0 s\,\nabla f(x_0))$. By the
fundamental theorem of calculus and (A3), for $s\le 2a/\|\nabla f(x_0)\|^2$,
$\varphi'(s)=-\langle\nabla f(x_0-\sigma_0 s\nabla f),\nabla f(x_0)\rangle\le-\|\nabla f(x_0)\|^2+\tfrac{C_g}{3}\|\nabla f(x_0)\|\le-\tfrac23\|\nabla f(x_0)\|^2$,
using (A2). Integrating to $s^*=2a/\|\nabla f(x_0)\|^2$ gives
$\varphi(s^*)\le\varphi(0)-\tfrac23\|\nabla f(x_0)\|^2 s^*=|f(x_0)|-\tfrac43 a\le a-\tfrac43a=-\tfrac13a<0$,
so $\sigma_0 f(x_1)<0$; the step length is $\|x_1-x_0\|=2a/\|\nabla f(x_0)\|\le 2a/C_g$.
\end{proof}

The Taylor and dual-norm refinements of \cite{advExt} sharpen the constant $6$ implicit in the
compatibility $a\le C_g^2/(6H)$ to $4$ or $2$ and give an $\ell_\infty$ signed-gradient
(FGSM) certificate; we use the plain version and note the refinements where relevant.

\subsection{The master theorem}
\begin{theorem}[Universal fragility of wide Gaussian networks]\label{thm:master}
Let $F$ be a network as in \Cref{sec:networks} --- fully connected with arbitrary biases,
convolutional, a fixed-depth residual network, or a simplified vision transformer (\Cref{sec:arch}) ---
with independent Gaussian weights in the signal-preserving normalization, a random readout
$w\sim N(0,I_{n_L}/n_L)$, and any $[0,1]$-slope activation satisfying the nondegeneracy of
\Cref{prop:gradlb} under \Cref{ass:nondeg}. Assume the wide regime $dL\log\bar n=o(\min_\ell n_\ell)$.
Then there is a constant $C=C(L,c_0,c_1,c)$ such that, with high probability over the weights,
\emph{simultaneously for every input $x$ on the sphere of radius $R$} with $F(x)\ne0$, the single
gradient step of \Cref{prop:engine} applied to $F$ changes the sign of $F$, with Euclidean length
\[
  \|x_1-x\|\ \le\ C\,\|x\|\,\sqrt{\frac{d\log\bar n}{n_L}}.
\]
For any fixed input $x$ the length improves to $C\,\|x\|/\sqrt{n_L}$ (no $d\log\bar n$ factor).
\end{theorem}

\begin{proof}
Fix the good event, of probability $1-o(1)$, on which the conclusions of \Cref{prop:readout,prop:gradlb,thm:stability}
all hold uniformly over the sphere. On this event, for every $x$:
\begin{itemize}
\item[(A1)] $|F(x)|\le a:=C_1\|x\|\sqrt{(d\log\bar n)/n_L}$ by \Cref{prop:readout};
\item[(A2)] $\|\nabla F(x)\|\ge C_g:=c_g\|w\|$ by \Cref{prop:gradlb}, and $\|w\|=1+o(1)$ w.h.p.;
\item[(A3)] for $\|x'-x\|\le 2a/C_g$, \Cref{thm:stability} gives
$\|\nabla F(x')-\nabla F(x)\|\le C\|w\|(\eta\log(1/\eta))^{1/2}$ with
$\eta\asymp(2a/C_g)\sqrt{\bar n}/\|x\|\asymp\sqrt{(d\log\bar n)\,\bar n/n_L}\,/\,c_g$. In the wide
regime this is $o(\|w\|)=o(C_g)$, so (A3) holds with room to spare once $\bar n\,d\log\bar n=o(n_L\,c_g^{2})$,
which is the wide regime for the readout width; for square-width networks $\bar n\asymp n_L$ we use the
sharpened stability radius, see \Cref{rem:widthcond}.
\end{itemize}
\Cref{prop:engine} then applies at every $x$, and the step length is $\|x_1-x\|\le 2a/C_g=
C\|x\|\sqrt{(d\log\bar n)/n_L}$. For a single fixed $x$ the net is unnecessary and \Cref{prop:readout}
gives $a=C_1\|x\|/\sqrt{n_L}$ directly, so the length is $C\|x\|/\sqrt{n_L}$.
\end{proof}

\begin{remark}[Width condition]\label{rem:widthcond}
The clean statement asks $\bar n\,d\log\bar n=o(n_L)$ for the stability step to dominate uniformly.
When all widths are comparable, $\bar n\asymp n_L\asymp n$, the naive form of (A3) is borderline; the
Taylor-remainder version of the engine (\cite[Thm 4.6]{advExt}, compatibility constant $2$ and only the
attack-line curvature) needs only $\|\nabla F(x')-\nabla F(x)\|=o(\|\nabla F\|)$ \emph{along the
segment}, which the directional form of \Cref{thm:stability} gives whenever $d\log n=o(n)$. So the
theorem holds in the standard wide regime $d\log n=o(n)$; we state the operator-norm form for
transparency.
\end{remark}

\begin{remark}[The scale]\label{rem:scale}
For a fixed input the radius $C\|x\|/\sqrt{n_L}$ matches the $\tilde O(\|x\|/\sqrt d)$ of
\cite{danielyS20} when the readout width tracks the input dimension, and the
$c\sqrt{\log(1/\delta)/d}$ of \cite{bartlettBC21}, but with three differences: the activation is any
$[0,1]$-slope function rather than ReLU only; the biases are arbitrary; and the constant in the depth
dependence is controlled by \Cref{prop:gradlb} rather than left exponential. The uniform-over-inputs
version, with its extra $\sqrt{d\log\bar n}$, is the price of ``every $x$ at once'' and is what the
arrangement bound of \Cref{lem:arrangement} makes affordable.
\end{remark}

\subsection{Multiclass}
For a multiclass network $F(x)=\arg\max_j\langle w_j,h_L(x)\rangle$ the attack targets the margin
against a fixed competitor.

\begin{theorem}[Multiclass fragility]\label{thm:multiclass}
Under the hypotheses of \Cref{thm:master} with $K$ classes and independent random readouts
$w_j\sim N(0,I/n_L)$, w.h.p.\ simultaneously for every $x$ on the sphere with predicted class
$j=F(x)$ there is a competitor $c\ne j$ and a single gradient step on the margin
$m_{j,c}(x)=\langle w_j-w_c,h_L(x)\rangle$ of length
$\|x_1-x\|\le C\|x\|\sqrt{(d\log(Kn))/n_L}$ with $F(x_1)\ne j$.
\end{theorem}

\begin{proof}
The margin $m_{j,c}$ is a network of the same class with readout $w_j-w_c\sim N(0,2I/n_L)$, so
\Cref{prop:readout} bounds $|m_{j,c(x)}(x)|$ for the runner-up $c(x)$ (the class with the largest
$\langle w_c,h_L\rangle$ below $j$) --- indeed for the runner-up the margin is the smallest positive
one, hence $\le C\|x\|\sqrt{(d\log(Kn))/n_L}$ by a union bound over the $K$ competitors --- while the
gradient lower bound and stability apply to $m_{j,c}$ verbatim. \Cref{prop:engine} on $m_{j,c}$ then
crosses $m_{j,c}=0$, at which point $c$ overtakes $j$. Uniformity is by the same net, with the extra
$\log K$ absorbed.
\end{proof}

\subsection{Dual norms}
The $\ell_2$ step is not the only one: the same three ingredients give an $\ell_\infty$ (signed
gradient) attack, which is the threat model of most empirical work.

\begin{corollary}[$\ell_\infty$ fragility]\label{cor:linf}
Under the hypotheses of \Cref{thm:master}, w.h.p.\ simultaneously for every $x$ the signed-gradient
step $x_1=x-\sgn(F(x))\,\tfrac{2a}{\|\nabla F(x)\|_1}\sgn(\nabla F(x))$ changes the sign of $F$, with
$\|x_1-x\|_\infty\le 2a/\|\nabla F(x)\|_1$. Since $\|\nabla F(x)\|_1\ge\|\nabla F(x)\|_2\ge c_g\|w\|$
and typically $\|\nabla F\|_1\asymp\sqrt d\,\|\nabla F\|_2$, the $\ell_\infty$ radius is
$O\!\bigl(\|x\|\sqrt{\log\bar n}/(\sqrt{n_L}\,)\bigr)$ per coordinate.
\end{corollary}

\begin{proof}
Apply the dual-norm engine \cite[Thm 4.7, Cor 4.8]{advExt} with $\|\cdot\|_{\mathcal X}=\ell_\infty$,
dual $\ell_1$, curvature bounded along the signed-gradient segment by \Cref{thm:stability}; the
compatibility $2a\,H_\infty<\|\nabla F\|_1^2$ holds because $a=o(1)$ while $\|\nabla F\|_1\ge c_g$ and
$H_\infty=o(1)$ in the wide regime.
\end{proof}

\section{Certificates and the certification gap}\label{sec:certificates}

Fragility has a robustness companion. The same operator-norm bracket that upper-bounds the network
Lipschitz constant gives, at every input, a radius inside which no perturbation can change the
output; this is a genuine certificate, valid for all inputs with no exceptional set. Comparing it with
the fragility radius of \Cref{thm:master} shows that the $1/\sqrt n$ scale is sharp in both
directions. We then prove the theorem that explains why one cannot do better by enumerating gate
patterns: any certificate required to hold over all patterns of a fixed density inherits the
worst-case conditioning, which is exponentially small by the product rate of \Cref{thm:product}.

\subsection{The Lipschitz certificate}
\begin{theorem}[Uniform robustness certificate]\label{thm:certificate}
Under the hypotheses of \Cref{thm:master} (fragility not required; only the Gaussian weights and the
wide regime), let
\[
  L_{\mathrm{net}}=\|w\|\prod_{\ell=1}^L\bigl(2\sqrt{\kappa(c_1)}\bigr)
\]
with $c_1$ the gate-density upper bound of \Cref{ass:nondeg}. Then w.h.p.\ $F$ is
$L_{\mathrm{net}}$-Lipschitz on the sphere, and consequently, for every input $x$, no perturbation of
Euclidean norm less than $|F(x)|/L_{\mathrm{net}}$ changes the sign of $F$. This certified radius holds
for \emph{all} inputs with no exceptional set.
\end{theorem}

\begin{proof}
For any $x,x'$ the difference $F(x')-F(x)$ is bounded by the operator norm of the Jacobian integrated
along the segment; on each cell of the arrangement (\Cref{lem:arrangement}) the Jacobian is a fixed
product of gated blocks, whose operator norm is at most
$\|w\|\prod_\ell\max_{|T_\ell|,|S_\ell|\le c_1 n_\ell}\|P_{T_\ell}W_\ell P_{S_\ell}\|_\op\le
\|w\|\prod_\ell(2\sqrt{\kappa(c_1)}+o(1))$ by \Cref{prop:A2}, uniformly over cells by the union bound
(the count is $(C\bar n)^{dL}$ and each operator-norm bound fails with probability $e^{-cn}$). Since
the gate densities are at most $c_1$ everywhere by \Cref{ass:nondeg}, this bounds the Lipschitz
constant on every cell, hence on the sphere. A perturbation smaller than $|F(x)|/L_{\mathrm{net}}$
changes $F$ by less than $|F(x)|$, so cannot change its sign.
\end{proof}

The subgaussian version, with an absolute constant $C(K)$ in place of $2\sqrt{\kappa(c_1)}$, holds for
any i.i.d.\ subgaussian weights of parameter $K$, by the subgaussian operator-norm bracket of
\Cref{app:brackets}; this is what makes the certificate meaningful for non-Gaussian and lightly-trained
weights.

\begin{corollary}[Sharpness of the $1/\sqrt n$ scale]\label{cor:sharp}
Under the hypotheses of \Cref{thm:master} with a random readout, for a fixed input $x$ (equivalently, for
all but a vanishing fraction of the sphere, excluding a neighborhood of the decision boundary where
$|F(x)|$ is atypically small), w.h.p.\ the certified radius and the fragility radius bracket each other:
\[
  \frac{c\,\|x\|}{\sqrt{n_L}\,L_{\mathrm{net}}}\ \le\ \bigl(\text{largest certified radius at }x\bigr)\ \le\
  \bigl(\text{attack radius at }x\bigr)\ \le\ \frac{C\,\|x\|\,\sqrt{d\log\bar n}}{\sqrt{n_L}} .
\]
Both scale as $\|x\|/\sqrt{n_L}$ up to the constants $L_{\mathrm{net}}=(2\sqrt{\kappa})^L$ and
$\prod\sqrt{\kcheck}$; the gap between them is $L_{\mathrm{net}}\big/\prod_\ell\sqrt{\kcheck(\alpha_\ell)}$,
a product of the two-sided gate constants and no power of $n$.
\end{corollary}

\begin{proof}
The lower bound is \Cref{thm:certificate} with $|F(x)|\ge c\|x\|/\sqrt{n_L}$, which holds w.h.p.\ at a
fixed $x$ because $F(x)\sim N(0,\|h_L\|^2/n_L)$ is not atypically small (anticoncentration of a
Gaussian); the upper bound is \Cref{thm:master}. The ratio of the two radii is
$L_{\mathrm{net}}\cdot(\text{attack constant})=(2\sqrt\kappa)^L\prod\sqrt{\kcheck}^{-1}$ up to
logarithms, independent of $n$.
\end{proof}

\subsection{The certification gap}
The certificate of \Cref{thm:certificate} is loose by the factor
$L_{\mathrm{net}}/\prod\sqrt{\kcheck}$ because it uses the worst-case operator norm at every step. One
might hope to sharpen it by certifying each gate pattern separately, using the actual conditioning of
the realized Jacobian rather than the worst pattern's. The product rate shows this cannot help beyond a
point: a certificate valid over all density-$\alpha$ patterns must survive the worst-conditioned one,
and that one is exponentially singular.

\begin{theorem}[Certification gap]\label{thm:cert-gap}
Consider certifying a depth-$k$ gated linear network $x\mapsto D_kW_k\cdots D_1W_1x$ with Gaussian
weights, by any method whose guarantee must hold uniformly over all gate tuples
$(D_1,\dots,D_k)$ of density $\alpha$. The smallest conditioning such a method must tolerate is the
least singular value of the worst pattern, $M_N^{(k)}(\alpha)$, which by \Cref{thm:product} is at most
$e^{-2H(\alpha)N/\dF+o(N)}$ w.h.p., and equal to $e^{-(k+1)H(\alpha)N/\dF+o(N)}$ at the sharp
(numerically determined) product rate. Consequently any pattern-uniform margin- or inversion-based
certificate degrades relative to the typical-pattern conditioning $N^{-(k+1)/2}$ by a factor at least
$e^{2H(\alpha)N/\dF+o(N)}$, and by $e^{(k+1)H(\alpha)N/\dF+o(N)}$ at the sharp rate: the union that a
uniform certificate performs is dominated by the single worst pattern, whose conditioning sets the rate.
\end{theorem}

\begin{proof}
A certificate uniform over the $\binom N{\alpha N}^{k+1}=e^{(k+1)H(\alpha)N+o(N)}$ patterns must in
particular be valid for the pattern minimizing $\smin$, whose value is $M_N^{(k)}(\alpha)$; by
\Cref{thm:product} this is $\le e^{-2H(\alpha)N/\dF+o(N)}$ w.h.p.\ (the proved achievability of
\Cref{prop:prod-lower}), and $=e^{-(k+1)H(\alpha)N/\dF+o(N)}$ at the sharp rate. A margin certificate
scales the certified radius by the least singular value of the map it inverts, so its worst-case radius
is smaller than the typical-pattern radius (governed by the fixed-support hard edge $N^{-(k+1)/2}$,
\Cref{prop:A4} and its least-sv companion) by the ratio $N^{-(k+1)/2}/M_N^{(k)}$, which is at least
$e^{2H(\alpha)N/\dF+o(N)}$ and equals $e^{(k+1)H(\alpha)N/\dF+o(N)}$ at the sharp rate. Either way it is
exponential in $N$.
\end{proof}

\begin{remark}[Interpretation]\label{rem:cert-interp}
\Cref{thm:cert-gap} is the precise sense in which the adversarial and robust pictures are two sides of
one random-matrix law. Fragility is the statement that the \emph{realized} pattern already conditions
the network badly enough for a gradient step to work; the certification gap is the statement that no
pattern-agnostic certificate can be sharp, because the \emph{worst} pattern is exponentially worse than
the realized one, at a rate that is at least $2H(\alpha)/\dF$ unconditionally and $(k+1)H(\alpha)/\dF$ at
the sharp product rate of \Cref{thm:product}. The only way
to beat the gap is to certify at the realized pattern --- i.e.\ to use the actual input, which is what a
per-input Lipschitz-of-the-realized-Jacobian certificate does; that recovers the $1/\sqrt n$ scale of
\Cref{cor:sharp} but is not uniform over patterns, and hence not robust to a gate flip, which is
exactly what an adversary induces.
\end{remark}

\subsection{Certified global robustness of sparse gates}
When the gate density is small enough that even the optimized product norm contracts, the network is
robust everywhere, with no small-perturbation attack possible.

\begin{theorem}[Sparse-gate robustness]\label{thm:sparse-robust}
Suppose \Cref{ass:nondeg} holds with $c_1<\alpha^*=0.042672$ (the root of $\kappa(\alpha^*)=1/4$). Then
w.h.p.\ the network Lipschitz constant $L_{\mathrm{net}}=(2\sqrt{\kappa(c_1)})^L\|w\|$ tends to $0$ as
$L\to\infty$, and for $L\ge L_0(c_1)$ the network computes a function whose oscillation on the sphere is
$o(1)$: no input has a certified radius smaller than $\Omega(1)$ relative to the output scale, so the
$1/\sqrt n$ fragility of \Cref{thm:master} is absent. Concretely, every perturbation of norm below
$|F(x)|/(2\sqrt{\kappa(c_1)})^L$ is harmless, and this radius grows with depth.
\end{theorem}

\begin{proof}
By \Cref{prop:contract}, $\kappa(c_1)<1/4$ makes $(2\sqrt{\kappa(c_1)})^L\to0$, so
$L_{\mathrm{net}}\to0$ and \Cref{thm:certificate} gives a certified radius $|F(x)|/L_{\mathrm{net}}$
that grows geometrically in $L$. The oscillation bound is the same operator-norm estimate applied to
$F(x)-F(x')$ over the whole sphere.
\end{proof}

This is the positive counterpart of the fragility theorems: sparsity of activation, quantified by the
gate density falling below $\alpha^*$, certifies global robustness, and the threshold is exactly the
contraction threshold of the optimized submatrix product. It is consistent with the converse of
\cite{bartlettBC21}, that networks deep enough to be near-constant are not fragile, and sharpens it to a
density condition with an explicit constant.

\section{Architectures}\label{sec:arch}

The master theorem was stated for a generic feedforward network; here we verify the two structural
hypotheses it uses --- the uniform gain bracket and the gradient lower bound with stability --- for
convolutional, residual, and transformer networks. In each case the point is that the input--output
Jacobian is again a product of gated Gaussian blocks, possibly structured, to which the estimates of
\Cref{sec:gains} apply, and that the biases and the choice of $[0,1]$-slope activation change only the
gate densities.

\subsection{Convolutional networks}
A convolutional layer applies a Gaussian filter bank followed by a gate. With $c$ channels and spatial
size $s$ the layer acts on $\R^{cs}$ by a block matrix whose blocks are (multichannel) circulant or
Toeplitz Gaussian matrices; the gate keeps the active pixel-channels. The Jacobian is therefore a
product of gated structured-Gaussian blocks.

\begin{theorem}[Convolutional fragility]\label{thm:cnn}
Let $F$ be a convolutional network with independent Gaussian filters, arbitrary biases, and any
$[0,1]$-slope activation, in the wide regime $d\log(cs)=o(cs)$ (effective width $cs$). Then the
conclusions of \Cref{thm:master} and \Cref{thm:certificate} hold: w.h.p.\ simultaneously for every
input a single gradient step of length $O(\|x\|\sqrt{(d\log cs)/(cs)})$ crosses the boundary, and every
input has a certified radius of the same order.
\end{theorem}

\begin{proof}[Proof sketch]
Two ingredients change relative to the fully connected case. First, the operator-norm bracket of
\Cref{prop:A2} must be applied to structured Gaussian blocks; the upper bound survives because the
circulant/Toeplitz structure only decreases the operator norm relative to a full Gaussian of the same
variance (the filter has $O(1)$ taps), and the lower bound survives because a fixed input direction is
mapped to a Gaussian vector whose top/bottom pixel-channels carry $\kappa,\kcheck$ of its mass as
before. Second, the gradient lower bound uses the nondegeneracy recursion of \Cref{prop:gradlb}
channelwise; this is exactly the content of \cite{danielyM26} for $C^2$ activations, and the
$[0,1]$-slope extension is by the convex-hull reduction of \Cref{lem:hull}. We do not use any Fourier
diagonalization of the nonlinear network --- only operator-norm bounds, which the frequency structure of
convolutions makes no worse --- so the barrier of \cite[Rem 10.2]{advExt} does not obstruct the argument.
The stability lemma \Cref{thm:stability} applies verbatim, with the anticoncentration of the
pixel-channel preactivations replacing that of the fully connected preactivations.
\end{proof}

\subsection{Residual networks}
A fixed-depth residual network has blocks $h_\ell=h_{\ell-1}+t^{-1/2}\sigma(\Lambda_\ell h_{\ell-1})$
with Gaussian $\Lambda_\ell$ and the standard residual scaling. Its block Jacobian is
$A_\ell=I+t^{-1/2}D_\ell\Lambda_\ell$, a gated near-identity, and the product $\prod_\ell A_\ell$ has
least and largest singular values pinched around $1$.

\begin{theorem}[Residual fragility, ReLU and general slopes]\label{thm:resnet}
Let $F$ be a fixed-depth residual network with $t$ blocks, Gaussian $\Lambda_\ell$ with $\|\Lambda_\ell\|_\op\le M$,
any $[0,1]$-slope activation, and arbitrary biases, with $t>S_1^2M^2$ (so the near-identity is
invertible), where $S_1=\|\sigma'\|_\infty\le1$. Set $\nu=1-S_1M/\sqrt t$, $\alpha=1+S_1M/\sqrt t$. Then
w.h.p.\ $\smin(J_h)\ge\nu^t$ and $\|J_h\|_\op\le\alpha^t$ uniformly over inputs, and the conclusions of
\Cref{thm:master,thm:certificate} hold with $C_g=\nu^t\|w\|$ and $L_{\mathrm{net}}=\alpha^t\|w\|$.
\end{theorem}

\begin{proof}[Proof sketch]
The deterministic block bounds $\smin(A_\ell)\ge\nu$, $\|A_\ell\|_\op\le\alpha$ are \cite[Prop 5.6]{advExt};
they hold for any $[0,1]$-slope $\sigma$ because they use only $\|D_\ell\|_\op\le S_1\le1$ and the gate
support, not the sign structure. What \cite{advExt} could not do for ReLU was the gradient lower bound,
which it obtained only in expectation and only for odd $\sigma$ via a cross-term cancellation. Here the
lower bound $\smin(J_h)\ge\nu^t$ is deterministic given the weights, so it holds for ReLU and every
$[0,1]$-slope activation with no oddness: the near-identity structure makes each block invertible with
$\smin\ge\nu>0$ regardless of the gate. The stability bound uses the same telescoping as
\Cref{thm:stability}, with the residual scaling $t^{-1/2}$ making the per-block Hessian
$H_{\mathrm{res}}=S_2M^2\sqrt t\,\alpha^{2t-2}$ (\cite[Prop 5.6(c)]{advExt}), and the compatibility
$a\le\nu^{2t}/(6H_{\mathrm{res}}\alpha^t)$ holds in the wide regime because $a=O(\|x\|/\sqrt{n_L})\to0$.
\end{proof}

The residual case is the one where the paper most directly strengthens \cite{advExt}: it removes the
oddness and bias-free restrictions and covers ReLU, precisely because the gradient lower bound is now
deterministic (near-identity invertibility) rather than an odd-activation expectation.

\subsection{Vision transformers}
A simplified vision transformer, in the sense of \cite[\S5.3]{advExt} (bias-free linear maps, no
positional embedding, affine-free normalization, odd MLP activation), has a Jacobian that is a product
of near-identity residual blocks whose linear part is a gated value-path plus a low-rank attention
correction. The attention weights are a softmax of a bilinear form in the input, and the difficulty
that \cite{advExt} left open (Remark 5.20 there) was that at standard temperature the softmax logits are
order one, so the attention Jacobian does not concentrate near a fixed linear map. In the wide regime
the value-path constants do concentrate, which lets us discharge the compact-set hypotheses of their
conditional theorem at random initialization.

\begin{theorem}[Simplified-ViT fragility at random init, wide regime]\label{thm:vit}
Let $F$ be a simplified vision transformer with $L$ layers, $h$ heads, key dimension $d_k$, token count
$n_{\mathrm{tok}}$, and Gaussian (Xavier) weights, in the wide regime $hd_k\to\infty$ with
$d\log(hd_k)=o(hd_k)$. Then w.h.p.\ the value-path operator norms $q_\ell,k_\ell,m_\ell$ of
\cite[Thm 5.16]{advExt} are $O(1)$, the standard-temperature Jacobian and Hessian bounds
$\Gamma_\ell,\Xi_\ell$ there are finite and $O(1)$ uniformly over the Frobenius ball of radius
$R\asymp\sqrt{n_{\mathrm{tok}}d}$, and on the local tube of \cite[Rem 5.17]{advExt} generated by the
attack neighborhoods the conclusions of \Cref{thm:master,thm:certificate} hold: a single gradient step of
Frobenius length $O(R/\sqrt{hd_k})$ crosses the boundary.
\end{theorem}

\begin{proof}[Proof sketch]
The constants $q_\ell=\|W_{Q,\ell}\|_\op$, $k_\ell=\|W_{K,\ell}\|_\op$, $m_\ell=\|W_{V,\ell}W_{O,\ell}\|_\op$
of \cite[Thm 5.16]{advExt} are operator norms of Gaussian matrices, hence $O(1)$ w.h.p.\ by the
Bai--Yin edge; substituting into their $\Gamma_\ell=m_\ell(\sqrt{n_{\mathrm{tok}}}+RL_{P,\ell})$,
$\Xi_\ell=m_\ell(RH_{P,\ell}+2L_{P,\ell})$ gives finite $O(1)$ values on $R\asymp\sqrt{n_{\mathrm{tok}}d}$
because $L_{P,\ell},H_{P,\ell}=O(R/\sqrt{d_k})$ times the universal softmax constants. The
near-identity structure ($\nu_{\mathrm{std}}=1-\beta\Gamma>0$ with $\beta=L^{-1/2}$) then gives the
gradient lower bound $\|\nabla F\|\ge\nu_{\mathrm{std}}^L\|w\|$ as in the residual case, deterministically
given the weights. What remains is the forward invariance of a set containing the attack trajectories;
we replace the global compact-set hypothesis of \cite[Thm 5.16]{advExt} by the local tube of their
Remark 5.17, whose existence follows from the stability estimate \Cref{thm:stability} applied to the
value-path (the attention correction is $O(\beta)$ and does not move the trajectory out of an
$O(R/\sqrt{hd_k})$ tube). This resolves the wide-regime, simplified, standard-temperature case of their
open problem; the fully standard transformer (biases, positional embeddings, affine normalization,
GELU) remains open, as those break the balance that supplies the small output.
\end{proof}

\begin{remark}
\Cref{thm:vit} is a partial resolution of \cite[Rem 5.20]{advExt}: it shows the standard-temperature
constants are $O(1)$ w.h.p.\ at random init and that forward invariance can be localized to a tube, so
the conditional theorem becomes an unconditional wide-regime one for the simplified architecture. The
random readout route of \Cref{prop:readout} additionally removes the oddness hypothesis on the MLP that
their balance argument needed for the small output, at the cost of the readout being random.
\end{remark}

\section{Depth: phase transitions and the dynamical limit}\label{sec:depth}

The gate density controls a phase transition in depth. Above a density threshold the optimized product
norm can grow and fragility persists to any depth; below it, the norm contracts and the network becomes
globally robust and eventually near-constant. The threshold is the contraction constant of the
optimized submatrix product, and the transition is sharp in the density. We also record the
weight-tied limit, where one Gaussian matrix is reused across depth and the long-time behaviour is set
by the eigenvalue radius of cyclic submatrix products rather than the operator norm.

\subsection{Density thresholds}
\begin{theorem}[Depth phase transition]\label{thm:phase}
Consider a depth-$L$ gated Gaussian network with all gate densities equal to $\alpha$, in the
signal-preserving normalization. Let $\alpha^*=0.042672$ be the root of $\kappa(\alpha)=1/4$ and, for
fixed support, $\alpha^*_k=(k^k/(k+1)^{k+1})^{1/k}$.
\begin{enumerate}
\item If $\alpha<\alpha^*$, then w.h.p.\ the optimized Jacobian operator norm $\prod$-bracket
contracts, $\max_{S_\bullet}\|J\|_\op\le(2\sqrt{\kappa(\alpha)})^L\to0$, and by \Cref{thm:sparse-robust}
the network is globally robust: every input has certified radius growing with depth, and no
single-step attack succeeds for large $L$.
\item If $\alpha^*<\alpha<1$, the operator-norm bracket does not force contraction: the optimized
product norm can exceed one, so fragility of the type in \Cref{thm:master} is possible at every depth,
and the realized-pattern Jacobian is $\Theta(1)$ so the fragility radius stays $\Theta(\|x\|/\sqrt n)$.
\end{enumerate}
\end{theorem}

\begin{proof}
Part (1) is \Cref{prop:contract} and \Cref{thm:sparse-robust}. For part (2), at $\alpha>\alpha^*$ the
upper bracket $(2\sqrt{\kappa(\alpha)})^L$ exceeds one and grows, so it does not certify contraction;
the realized-pattern Jacobian has the signal-preserving $\Theta(1)$ norm of \Cref{prop:gradlb}, so the
gradient lower bound and hence \Cref{thm:master} apply at every fixed depth.
\end{proof}

The gap between $\alpha^*=0.043$ (optimized) and $\alpha^*_k$ (fixed support, e.g.\ $0.535$ at $k=4$)
is the entropy the adversary buys by choosing patterns: a fixed half-support product of length $4$
already contracts, while the optimized-support threshold stays near $\alpha^*$, because the optimizer
can find the rare non-contracting pattern. This is the operator-norm face of the certification gap.

\subsection{The weight-tied dynamical limit}
When one Gaussian matrix $W$ is reused at every layer, $x_{t+1}=\sigma(Wx_t)$, the per-step map is a
gated copy of $W$ and the long-time behaviour is a question about cyclic products of its submatrix
blocks. The relevant facts, from the companion analysis \cite{relumap}, are that the state norm
contracts at the fresh-Gaussian rate, that the activation density is $1/2$, and that no expanding
direction survives:

\begin{proposition}[Weight-tied contraction; {\cite{relumap}}]\label{prop:tied}
Let $W$ be $N\times N$ with i.i.d.\ $N(0,1/N)$ entries and $x_{t+1}=\sigma(Wx_t)$ for ReLU $\sigma$.
For every fixed horizon $T$ and deterministic input, w.h.p.\ $\|x_t\|^2\to 2^{-t}$ for $t\le T$ and the
gate density is $1/2$. For every fixed $r>2^{-1/2}$, w.h.p.\ there is no ReLU ray of multiplier
$\ge r$, so every trajectory whose normalized direction converges decays at every rate exceeding
$2^{-1/2}$. The sustained gain of a period-$p$ cycle is the eigenvalue radius of the cyclic submatrix
product, $2^{-p/2}$ for asymptotically free blocks by Haagerup--Larsen, below one for all $p\ge2$.
\end{proposition}

The relevance to fragility is that the output at depth $t$ is $|F|\asymp\|x_t\|\asymp2^{-t/2}\|x_0\|$
(random readout) while the Jacobian also contracts, so the ratio --- the fragility radius --- is set by
the competition between the two, and the eigenvalue-radius contraction of \Cref{prop:tied} shows the
Jacobian never expands along a cycle even though its submatrix operator norm exceeds one. The
weight-tied network is thus in the robust phase for the sustained (deep, cyclic) part of its dynamics,
and the fragility is a finite-horizon phenomenon governed by the transient submatrix products. This is
the dynamical-systems reading of the density threshold: the reused-weight map has effective density
$1/2>\alpha^*$, so it is fragile at finite depth, but its cyclic eigenvalue radius is contractive, so it
does not sustain an expanding adversarial direction.

\begin{remark}[Why the singular value and the eigenvalue disagree]\label{rem:sv-vs-ev}
The half-by-half submatrix operator norm is $\sqrt2>1$, and its optimized-support maximum is larger
still, up to $2\sqrt{\kappa(1/2)}=1.927$; yet the cyclic eigenvalue radius is $2^{-p/2}<1$. The gap is
the non-normality of the gated product, and it is why bounds on submatrix \emph{singular} values alone
cannot decide the infinite-time dynamics --- the ReLU sign constraint, which selects a specific
nonnegative-cone-feasible pattern, is what collapses the operator-norm bracket to the eigenvalue radius.
The same distinction underlies the certification gap: an adversary works with singular values (worst
pattern, operator norm), while the sustained dynamics works with eigenvalues (realized cycle, spectral
radius).
\end{remark}

\section{Trained and semi-random networks}\label{sec:trained}

The theorems so far assume Gaussian weights. Trained networks are not Gaussian, but empirically their
weight matrices are a Gaussian bulk plus a small number of outlier directions, and the balance route to
the small output does not need Gaussianity at all. This section records what survives training: a
perturbation-stability statement that degrades the Gaussian bounds gracefully under a bounded deviation
of the weights, the balance version of the fragility theorem for trained classifiers, and a discussion
of the numerical evidence.

\subsection{Perturbation stability}
\begin{theorem}[Stability under weight deviation]\label{thm:trained}
Let $W_\ell=W_\ell^0+E_\ell$ with $W_\ell^0$ Gaussian (satisfying the hypotheses of \Cref{thm:master})
and $E_\ell$ arbitrary with $\|E_\ell\|_\op\le\eps_\ell$. Then:
\begin{enumerate}
\item the Lipschitz certificate of \Cref{thm:certificate} holds with $L_{\mathrm{net}}$ replaced by
$\|w\|\prod_\ell(2\sqrt{\kappa(c_1)}+\eps_\ell)$;
\item the gradient lower bound of \Cref{prop:gradlb} survives, $\|\nabla F(x)\|\ge c_g\|w\|-\sum_\ell\eps_\ell G_\ell$,
where $G_\ell$ is the product of the operator-norm bounds of the other layers, provided the deviation
is small, $\sum_\ell\eps_\ell G_\ell<\tfrac12 c_g\|w\|$;
\item consequently the fragility conclusion of \Cref{thm:master} holds for the perturbed network under
this smallness condition, with the constants shifted accordingly.
\end{enumerate}
\end{theorem}

\begin{proof}
(1) Operator norms are subadditive, $\|D_\ell(W_\ell^0+E_\ell)P\|_\op\le\|D_\ell W_\ell^0P\|_\op+\eps_\ell$,
so each factor's bound increases by $\eps_\ell$ and the product bound is
$\prod_\ell(2\sqrt{\kappa(c_1)}+\eps_\ell)$. (2) Write the Jacobian of the perturbed network as the
Gaussian Jacobian plus a sum of $L$ terms, each replacing one $W_\ell^0$ by $E_\ell$ and bounding the
other factors by their operator norms; the correction is at most $\sum_\ell\eps_\ell G_\ell$ in norm,
which is subtracted from the Gaussian lower bound $c_g\|w\|$. (3) Under the smallness condition the
perturbed gradient is still $\ge\tfrac12c_g\|w\|$, and the small output and stability are perturbed by
the same $O(\sum\eps_\ell)$ amounts, so \Cref{prop:engine} applies with the shifted constants.
\end{proof}

The smallness condition is a low-rank/small-norm deviation from the Gaussian bulk, which is what the
Marchenko--Pastur analysis of trained weights \cite{shmaloRMT,shmaloAccuracy} identifies: the bulk of
the singular spectrum stays at the random edge while a few directions carry the learned signal, so
$\|E_\ell\|_\op$ is the top learned singular value, typically $O(1)$ and not growing with width, while
$c_g\|w\|$ and the certificate scale are $\Theta(1)$. The theorem then says training moves the
constants but not the $1/\sqrt n$ scale.

\subsection{The balance route for trained classifiers}
For a trained classifier the readout is not random, so the small output must come from balance. If the
trained classifier is accurate and its ground-truth labels are balanced, then its sign is balanced under
the input measure (pushed to the sphere), and \Cref{prop:balance} gives the small output for all but a
$\delta$-fraction of inputs. The gradient facts, however, need the weights to be near-Gaussian, which is
\Cref{thm:trained}. Combining:

\begin{corollary}[Trained-network fragility, informal]\label{cor:trained-frag}
Let $F$ be a trained classifier on the sphere whose weights are a Gaussian bulk plus a bounded deviation
(\Cref{thm:trained}) and whose sign is $p$-balanced. Then for all but a $\delta$-fraction of inputs a
single gradient step of length $O(\eps_{\mathrm{adv}}(p,\delta)\cdot L_{\mathrm{net}}/c_g)$ crosses the
boundary.
\end{corollary}

This is the trained-regime statement of \Cref{thm:master}: the small output is measure-theoretic (from
balance) rather than pointwise (from a random readout), and the gradient facts hold up to the weight
deviation. It is the theorem that the CIFAR experiments of \cite{advExt} were designed to probe.

\subsection{Numerical evidence from trained networks}
The experiments of \cite{advExt} train small CNNs, ResNets, and ViTs on CIFAR-10 and CIFAR-100 and
measure, along training, the one-step attack success, the predicted radius $2m/\|\nabla m\|$ against the
empirical attack radius, and the compatibility ratio. Three of their findings are exactly what the
present theory predicts and are worth restating in its language. First, the empirical attack radius is
monotone in the predicted radius $2m/\|\nabla m\|$ and stays below it, on both datasets --- the
conservative direction, since the prediction is an upper bound from the small-output estimate. Second,
the one-step vulnerability signatures appear early and persist after the networks become accurate, so
fragility is not a symptom of poor accuracy but a structural property of the map, as the balance-plus-%
gradient mechanism requires. Third, the transformer shows the strongest signatures despite the lowest
accuracy, consistent with the near-identity residual structure making its gradient lower bound most
robust. The experiments are diagnostics, not certified radii, and the frozen artifacts are catalogued in
\cite{advExt}; we use them as the trained-regime evidence for the scale that \Cref{thm:master} and
\Cref{cor:sharp} predict, and add the new numerical verification of the product rate in the next section.

\section{Numerical verification of the product rate}\label{sec:numerics}

The one new random-matrix constant in the paper --- the product least-singular-value rate
$(k+1)H(\alpha)/\dF$ of \Cref{thm:product} --- differs from the previously conjectured $2H(\alpha)/\dF$
starting at $k=2$, so it is worth checking directly. We do so by exhaustive minimization over all
support tuples at small $N$, where the exponent can be read from a single matrix, and by the cascade
construction at slightly larger $N$, where the achieving pattern can be built but exhaustive search is
infeasible.

\subsection{Setup}
For $N\in\{8,10\}$, density $\alpha=1/2$ (so $n=N/2$), and real Gaussian entries $N(0,1/N)$, we compute
the single-layer minimum $M_N^{(1)}=\min_{|S_0|=|S_1|=n}\smin(A_1[S_1,S_0])$ and the depth-two minimum
$M_N^{(2)}=\min_{S_0,S_1,S_2}\smin(A_2[S_2,S_1]A_1[S_1,S_0])$ over all support triples, reporting the
rates $\rho_k=-N^{-1}\log M_N^{(k)}$. For $N=8$ the triple minimization is over $70^3$ tuples and is
done exactly; for $N=10$ it is over $252^3$ tuples, so there we report the greedy cascade (the
two-sided-optimal first factor followed by the one-sided-optimal second factor), which is the
achieving construction of \Cref{app:product} and upper-bounds the exhaustive minimum. The discriminating
statistic is the ratio $\rho_2/\rho_1$: \Cref{thm:product} predicts $3H/2H=1.5$, the earlier conjecture
$2H/2H=1$.

\subsection{Results}
The measured rates are as follows (means over $8$ trials at $N=8$, $4$ at $N=10$; $2H(1/2)=1.386$,
$3H(1/2)=2.079$).

\begin{center}
\renewcommand{\arraystretch}{1.2}
\begin{tabular}{@{}lccc@{}}
\toprule
 & $\rho_1$ (exhaustive) & $\rho_2$ (exhaustive) & $\rho_2/\rho_1$ \\
\midrule
$N=8$  & $1.375$ & $2.087$ & $1.53$ \\
$N=10$ & $1.579$ & $2.338$ (cascade) & $1.50$ (cascade) \\
\bottomrule
\end{tabular}
\end{center}

At $N=8$, where the depth-two minimum is exact, $\rho_2=2.087$ sits on the predicted $3H=2.079$, and the
ratio $\rho_2/\rho_1=1.53$ is at the value $1.5$ that \Cref{thm:product} requires, cleanly separated from
the conjectured $1.0$. The greedy cascade tracks the exhaustive minimum from above (its rate is a lower
estimate of the true rate, and lies within finite-$N$ error), confirming that the joint-support optimization of \Cref{app:product} is the mechanism realizing the
minimum. At $N=10$ the cascade ratio is $1.50$. The single-layer rate $\rho_1$ fluctuates around
$2H=1.386$ at these small sizes (across seeds, means of $1.42$, $1.31$, $1.36$ at $N=8,10,12$), the
$o(N)$ correction to a first-passage exponent being of either sign; the ratio $\rho_2/\rho_1$, which
cancels the leading correction, is the robust quantity and confirms the linear-in-$k$ growth of the
exponent.

\subsection{The anisotropic cascade step}
The first moment of the cascade construction (\Cref{prop:aniso-mean}) predicts that the least singular
value of an anisotropic row ensemble $M$ with covariance $\Sigma$ of least eigenvalue $\lambda_{\min}$,
minimized over row supports, scales as $\sqrt{n\lambda_{\min}}$ times the isotropic rate. We check the
$\sqrt{\lambda_{\min}}$ prefactor directly. Taking $\Sigma=\diag(1,\dots,1,\eps^2)$ and computing the
exact minimum over row supports at $N=10$, $n=5$, the minimum is proportional to $\eps$ once $\eps$ is
small enough that the least eigenvalue dominates: the ratio $\min/(\eps\cdot\min|_{\eps=1})$ saturates to
a constant for $\eps\le 0.1$, and the log--log slope of $\min$ against $\eps$ is $1.0$--$1.15$ there,
matching the predicted linear dependence. This is the regime the cascade uses, where
$\lambda_{\min}=\smin(Q)^2$ is exponentially small.

The code is \texttt{numerics/verify\_product\_rate.py} and \texttt{verify\_anisotropic\_onesided.py};
both use only \texttt{numpy}. The verification refutes the product conjecture of \cite{shmaloProducts}
and supports \Cref{thm:product} and the first-moment reduction of \Cref{prop:aniso-mean}.

\section{Discussion}\label{sec:discussion}

The paper makes one claim in two registers. In random-matrix terms: the worst conditioning of a chain
of coordinate-optimized Gaussian submatrices decays at the sharp entropy rate $(k+1)H(\alpha)/\dF$, and
the optimized operator norm moves only to a bracket set by the retained-mass constants
$\kappa,\kcheck$. In network terms: a wide Gaussian network is fragile at the $1/\sqrt n$ scale at every
input, robust at exactly the reciprocal scale, and un-certifiable pattern-by-pattern below the same
entropy rate --- and all three are the same law read through the dictionary that gates are coordinate
selections and the Jacobian is a submatrix product.

Three features distinguish the account from its predecessors. It needs no oddness or bias-free
hypothesis: the small output comes either from a random readout, which is symmetry-free, or from
balance, which covers trained classifiers, and the gradient facts come from submatrix conditioning,
which does not see the sign structure. It is two-sided: the same bracket that produces the attack
produces a matching certificate, so the scale is pinned rather than bounded on one side. And it is
uniform over inputs, because the arrangement bound makes the union over realized gate patterns
affordable in the wide regime. The cost is the wide regime itself, $dL\log\bar n=o(\min_\ell n_\ell)$,
and the Gaussian (or subgaussian, or Gaussian-plus-bounded-deviation) weights; the fully trained,
proportional-width regime is reached only through the perturbation-stability theorem, which keeps the
scale but not the constants.

Several questions remain. The product rate is proved for equal densities and square bottlenecks; the
unequal-density cascade of \Cref{rem:unequal} should give the general rate, but the sharp constant for a
genuine rectangular bottleneck chain is open, as is the corresponding operator-norm edge for unequal
densities \cite{shmaloProducts}. The transformer theorem covers the simplified, standard-temperature
architecture at random init in the wide regime; the fully standard transformer, with biases, positional
embeddings, affine normalization, and GELU, breaks the balance that supplies the small output, and the
random-readout route does not obviously apply because the readout is a learned class token. The
certification gap is stated for pattern-uniform certificates; whether a certificate that adapts to the
realized pattern but is stable under a single gate flip can beat the $1/\sqrt n$ scale is the precise form
of the ``can we certify robustness'' question, and \Cref{rem:cert-interp} suggests the answer is no
without controlling the flip, which is what the adversary controls. Finally, the weight-tied dynamical
limit shows the eigenvalue radius, not the operator norm, governs the sustained gain; making the
infinite-time contraction unconditional requires the multi-time sign-decoupling estimate identified in
\cite{relumap}, which is the one genuinely open problem in the reused-weight theory.

The unifying point is methodological. Adversarial fragility has been explained by feature compression
\cite{fawzi18}, by computational hardness \cite{bubeckSS19}, and by isoperimetry
\cite{shafahi19,advExt}; the present account locates it in the spectral theory of the network's own
Jacobian, and shows that the same spectral object --- a product of Gaussian submatrix blocks, optimized
over supports --- controls the attack, the certificate, and the gap between them, quantitatively and with
matching constants. The random-matrix theorems were the missing ingredient; with them the geometry of
adversarial examples in deep networks is, in the wide Gaussian regime, a corollary of the conditioning
of coordinate-optimized Gaussian maps.


\appendix
\section{Proof of the two-sided single-layer rate}\label{app:twosided}

This appendix proves \Cref{thm:twosided}: for an $N\times N$ Gaussian matrix $A$ and
$M_N(\alpha)=\min_{|S|=|T|=n}\smin(A[T,S])$, $n=\lfloor\alpha N\rfloor$,
\[
  -\frac1N\log M_N(\alpha)\ \xrightarrow{\ \PP\ }\ \frac{2H(\alpha)}{\dF}.
\]
The upper bound is a first moment; the lower bound is a second moment over an alternately marked family
whose correlations are controlled by a sharp two-direction small-ball estimate. The argument is that of
\cite{shmaloSharp}, reproduced here so the product proof of \Cref{app:product} can quote its
construction and its \Cref{lem:twoslab}.

\subsection{First moment: the upper bound}
\begin{proposition}\label{prop:upper}
For every $\eps>0$, $\PP(M_N(\alpha)\le e^{-(2H(\alpha)/\dF+\eps)N})\to0$.
\end{proposition}
\begin{proof}
A single $n\times n$ Gaussian matrix $G$ has the hard-edge tail
\begin{equation}\label{eq:hardedge}
  \PP(\smin(G)\le t)\ \le\ C_\F(\sqrt n\,t)^{\dF},\qquad t>0,
\end{equation}
(real case Edelman and Rudelson--Vershynin; complex analogously, the least squared singular value having
bounded density near $0$ at scale $1/n$). Each $A[T,S]$ has the law of $G$, so a union bound over the
$\binom Nn^2=e^{2H(\alpha)N+o(N)}$ pairs gives $\PP(M_N(\alpha)\le t)\le e^{2H(\alpha)N+o(N)}t^{\dF}$;
at $t=e^{-(2H(\alpha)/\dF+\eps)N}$ the right side is $e^{-\dF\eps N+o(N)}\to0$.
\end{proof}

\subsection{The sharp two-direction small-ball estimate}
Let $g$ be a standard $\F$-Gaussian vector in $\F^D$ and $u,v$ unit vectors; write
$p(t)=\PP(|\zeta|\le t)$ for a standard $\F$-Gaussian scalar. The following is the one analytic
ingredient.

\begin{lemma}\label{lem:twoslab}
For all unit $u,v\in\F^D$ with $|\langle u,v\rangle|^2=R$, $s^2=1-R$, and all $0<t\le1$,
\begin{equation}\label{eq:cond}
  \PP\bigl(|\langle g,u\rangle|\le t,\ |\langle g,v\rangle|\le t\bigr)\ \le\ p(t)\,p(2t/s)\ \le\
  C_\F\,p(t)\,\min\bigl(1,(t/s)^{\dF}\bigr).
\end{equation}
If moreover $u,v$ are independent Haar-random unit vectors in $\F^D$, then averaging over their overlap,
for all $D\ge2$ and $0<t\le1$,
\[
  \PP\bigl(|\langle g,u\rangle|\le t,\ |\langle g,v\rangle|\le t\bigr)\ \le\ C_\F\,p(t)^2\,(1+\log(1/t)),
\]
with $C_\F$ independent of $D$.
\end{lemma}
\begin{proof}
Put $X=\langle g,u\rangle$, $Y=\langle g,v\rangle$; jointly $\F$-Gaussian, unit variances,
$\EE[Y\overline X]=\rho'$ with $|\rho'|^2=R$. Regress $Y=cX+sW$ with $|c|^2=R$ and $W$ a standard
$\F$-Gaussian independent of $X$. If $|X|\le t$ then $|cX|\le t$, so $|sW|=|Y-cX|\le|Y|+|cX|\le2t$ on the
event; hence the event is contained in $\{|X|\le t\}\cap\{|W|\le2t/s\}$, and independence of $X,W$ gives
$\PP\le p(t)p(2t/s)\le C_\F p(t)\min(1,(t/s)^{\dF})$ by \eqref{eq:pball}. For the averaged bound,
$R=|\langle u,v\rangle|^2\sim\Beta(\tfrac{\dF}2,\tfrac{(D-1)\dF}2)$, so $s^2=1-R\sim\Beta(\tfrac{(D-1)\dF}2,\tfrac{\dF}2)$,
and with $\alpha_0=\dF/2$, $b_0=(D-1)\dF/2$,
\[
  \EE[\min(1,(t/s)^{\dF})]=t^{\dF}\EE[s^{-\dF}\one_{\{s\ge t\}}]+\PP(s<t).
\]
The first expectation is $\tfrac1{B(b_0,\alpha_0)}\int_{t^2}^1 w^{b_0-\alpha_0-1}(1-w)^{\alpha_0-1}\dd w$;
at $D=2$ the exponent $b_0-\alpha_0-1=-1$ makes it $\le C_\F(1+\log(1/t))$, and at $D\ge3$ the exponent
exceeds $-1$, the integral converges, and the ratio of Beta functions is bounded uniformly in $D$
($\Gamma(x+a)/\Gamma(x)\asymp x^a$). The second is $\PP(\Beta(b_0,\alpha_0)<t^2)\le C_\F t^{2b_0}\le
C_\F t^{\dF}$ since $2b_0=(D-1)\dF\ge\dF$. Hence $\EE[\min(1,(t/s)^{\dF})]\le C_\F t^{\dF}(1+\log(1/t))$,
and $t^{\dF}\le p(t)/c_\F$ finishes.
\end{proof}

The averaging is packaged for reuse as follows; the factor $\Lambda$ carries polynomial overlap
prefactors.

\begin{lemma}\label{lem:averaged}
Let $j\ge1$, $q=\dF j$, and $\sigma^2\in[0,1]$ a random variable with, for some $\beta\ge\dF(j+1)/2$,
$\gamma\ge1$, $\Lambda\ge1$, $\PP(\sigma^2\le x)\le\Lambda(\gamma x)^\beta$ for all $x\in(0,1]$. Then for
$\sqrt\gamma\,t\le1$, $\EE[\min(1,(2t/\sigma)^q)]\le C_{\F,j}\Lambda\gamma^{q/2}t^q$.
\end{lemma}
\begin{proof}
Since $\gamma x\le1$ where the bound is nontrivial, take $\beta=\dF(j+1)/2$. Then
$\EE[\min(1,(2t/\sigma)^q)]\le\PP(\sigma^2\le4t^2)+(2t)^q\EE[\sigma^{-q}\one_{\{\sigma^2>4t^2\}}]$; the
first term is $\le\Lambda(4\gamma t^2)^\beta\le C\Lambda\gamma^{q/2}t^q$. Integrating the second by parts
and splitting at $w_*=(\gamma\Lambda^{1/\beta})^{-1}$, using $\beta-q/2=\dF/2$, both pieces are
$\le C\Lambda\gamma^{q/2}$; multiplying by $(2t)^q$ finishes.
\end{proof}

\subsection{The marked-row construction}
Fix $\eps\in(0,H(\alpha)/\dF)$, set $t=t_N=e^{-(2H(\alpha)/\dF-\eps)N}$, and a marked set
$M\subseteq[N]$ with $|M|=m_N=\lceil\sqrt N\rceil$. A \emph{configuration} is $c=(S,R,y)$ with $|S|=n$;
$y\in M$; $R\subseteq[N]\setminus M$, $|R|=n-1$. The rows $\{A_i|_S:i\in R\}$ span a hyperplane in
$\F^S\cong\F^n$ a.s.; let $u_c\in\F^S$ be a unit normal. Put
$E_c=\{|\langle A_y|_S,u_c\rangle|\le t\}$ and $W=\sum_c\one_{E_c}$.

\begin{lemma}\label{lem:suff}
If $W>0$ then $M_N(\alpha)\le t$.
\end{lemma}
\begin{proof}
On $E_c$, set $T=R\cup\{y\}$, $|T|=n$. For $i\in R$, $\langle A_i|_S,u_c\rangle=0$; hence
$\|A[T,S]u_c\|^2=|\langle A_y|_S,u_c\rangle|^2\le t^2$, so $\smin(A[T,S])\le t$.
\end{proof}

\begin{lemma}\label{lem:mean}
$\EE W=m_N\binom Nn\binom{N-m_N}{n-1}p(t)=e^{2H(\alpha)N+o(N)}t^{\dF}=e^{\dF\eps N+o(N)}\to\infty.$
\end{lemma}
\begin{proof}
Since $y\notin R$, $A_y$ is independent of $u_c$, so $\PP(E_c)=p(t)$; there are
$m_N\binom Nn\binom{N-m_N}{n-1}$ configurations, and $m_N=o(N)$ makes both binomials $e^{H(\alpha)N+o(N)}$.
\end{proof}

\subsection{Second moment and the lower bound}
\begin{proposition}\label{prop:lower}
For every $\eps>0$, $\PP(M_N(\alpha)>e^{-(2H(\alpha)/\dF-\eps)N})\to0$.
\end{proposition}
By \Cref{lem:suff} and Chebyshev, $\PP(M_N>t)\le\PP(W=0)\le\Var(W)/(\EE W)^2$, so it suffices that
$\Var(W)=o((\EE W)^2)$, i.e.\ $\EE[W^2]\le(\EE W)^2(1+o(1))$. Write $\mathcal G=\sigma(A_i:i\ne y)$ for a
pair with marked rows $y,y'$.

\emph{Distinct marked rows.} If $y\ne y'$, then $A_y$ enters neither $u_c$ (built from $R\not\ni y$), nor
$u_{c'}$ (from $R'\not\ni y$ since $y\in M$, $R'\subseteq[N]\setminus M$), nor $A_{y'}$; so $E_{c'}$ and
$u_c$ are $\mathcal G$-measurable while $\PP(E_c\mid\mathcal G)=p(t)$, giving $\PP(E_c\cap E_{c'})=p(t)^2$.
Summed, these pairs contribute at most $(\#\text{configs})^2p(t)^2=(\EE W)^2$.

\emph{Same marked row.} If $y=y'$, condition on $\mathcal G$; then $u_c,u_{c'}$ are determined and, with
$\tilde u_c,\tilde u_{c'}$ the zero-padded normals, \eqref{eq:cond} gives
$\PP(E_c\cap E_{c'}\mid\mathcal G)\le C_\F p(t)\min(1,(t/s)^{\dF})$, $s^2=1-R$,
$R=|\langle\tilde u_c,\tilde u_{c'}\rangle|^2$. Set $D_{\mathrm c}=n-|S\cap S'|$, $D_{\mathrm r}=n-|R\cap R'|$
(so $D_{\mathrm r}\ge1$). The overlap has two tail bounds.

\begin{lemma}\label{lem:overlaptail}
For all $x\in(0,1]$:
{\rm(i)} if $D_{\mathrm r}\ge2$, $\PP(R>1-x)\le C_\F x^{(D_{\mathrm r}-1)\dF/2}$, uniformly in
$D_{\mathrm r},N$;\ {\rm(ii)} if $D_{\mathrm c}\ge1$, $\PP(R>1-x)\le n(C_\F nx/D_{\mathrm c})^{D_{\mathrm c}\dF/2}$.
\end{lemma}
\begin{proof}
(ii) $R\le\|u_c|_{S\cap S'}\|^2=1-Z$, $Z=\|u_c|_{S\setminus S'}\|^2$; $u_c$ is Haar on the sphere of
$\F^S$, so $Z\sim\Beta(a,b)$, $a=D_{\mathrm c}\dF/2$, $b=(n-D_{\mathrm c})\dF/2$. Then
$\PP(Z\le x)\le x^a/(aB(a,b))\le n(C_\F nx/D_{\mathrm c})^a$ by $\Gamma(a+b)\le\Gamma(b)(a+b)^{a+1}$ and
$\Gamma(a+1)\ge(a/e)^a$. (i) Condition on the rows $R$ (determining $u_c$); the $D_{\mathrm r}-1$ rows in
$R'\setminus R$ project to independent Gaussians in $W_0=(\Span\{A_i|_{S'}:i\in R\cap R'\})^\perp$, of
dimension $D_{\mathrm r}$, and $u_{c'}$ is Haar on the sphere of $W_0$ orthogonal to them, i.e.\ Haar on
a $D_{\mathrm r}$-sphere; so $R\le|\langle\hat w,u_{c'}\rangle|^2=1-Z'$ with $Z'\sim\Beta(b_0,a_0)$,
$b_0=(D_{\mathrm r}-1)\dF/2$, $a_0=\dF/2$, giving $\PP(Z'\le x)\le C_\F x^{b_0}$.
\end{proof}

Now average $C_\F p(t)\min(1,(t/s)^{\dF})$ over three regions, using \Cref{lem:averaged} ($j=1$, $q=\dF$)
and the counting facts that the fraction of ordered pairs $(S,S')$ with $D_{\mathrm c}\le a$ is
$\le N^{2a}e^{-H(\alpha)N+o(N)}$, likewise for rows.

\emph{$D_{\mathrm r}\ge3$:} the row bound gives $\beta=(D_{\mathrm r}-1)\dF/2\ge\dF$, $\gamma=C_\F$,
$\Lambda=1$, so $\PP(E_c\cap E_{c'})\le C_\F p(t)t^{\dF}\le C_\F p(t)^2$; summed over $m_N$ marked rows,
$\le C_\F(\EE W)^2/m_N=o((\EE W)^2)$.

\emph{$D_{\mathrm r}\le2$, $D_{\mathrm c}\ge2$:} the column bound gives
$\PP(E_c\cap E_{c'})\le N^Cp(t)^2$; the rows are pinned ($\le N^4e^{-H(\alpha)N}$ fraction), so the
contribution is $N^Ce^{-H(\alpha)N+o(N)}(\EE W)^2=o((\EE W)^2)$.

\emph{$D_{\mathrm r}\le2$, $D_{\mathrm c}\le1$:} the crude bound $\PP(E_c\cap E_{c'}\mid\mathcal G)\le p(t)$
gives ratio $C_\F t^{-\dF}=C_\F e^{(2H(\alpha)-\dF\eps)N}$ to $p(t)^2$; both coordinates pinned
($\le N^Ce^{-2H(\alpha)N}$ fraction), so the contribution is $N^Ce^{-\dF\eps N+o(N)}(\EE W)^2=o((\EE W)^2)$,
including the diagonal.

Hence $\EE[W^2]\le(\EE W)^2(1+o(1))$, so $\PP(W=0)\to0$ and \Cref{prop:lower} follows.
\Cref{thm:twosided} is \Cref{prop:upper,prop:lower}. \qed

\begin{remark}\label{rem:marked}
The marked set forbids a configuration's normal-defining rows from the marked set, which removes the
$\Theta(1)$ fraction of genuinely correlated pairs in which each configuration's marked row is the
other's normal row, leaving only the same-marked-row pairs of weight $\Theta(1/m_N)$. Any $m_N\to\infty$,
$m_N=o(N)$ works.
\end{remark}

\section{The product least singular value}\label{app:product}

This appendix proves the bracket of \Cref{thm:product} and explains the sharp constant. Recall
$M_N^{(k)}(\alpha)=\min_{S_0,\dots,S_k}\smin(P_{S_k}A_k\cdots P_{S_1}A_1P_{S_0})$ with all supports of
size $n=\lfloor\alpha N\rfloor$, and set $\rho_k=-N^{-1}\log M_N^{(k)}(\alpha)$. \Cref{thm:product}
asserts $2H(\alpha)/\dF\le\liminf\rho_k\le\limsup\rho_k\le(k+1)H(\alpha)/\dF$, with the upper end the
sharp value. The two ends are \Cref{prop:prod-upper,prop:prod-lower}; the sharp value and the
mechanism are \Cref{rem:achieve}.

\subsection{The union bound: $\rho_k\le(k+1)H/\dF$}
\begin{proposition}\label{prop:prod-upper}
For every $\eps>0$, $\PP\bigl(M_N^{(k)}(\alpha)\le e^{-((k+1)H(\alpha)/\dF+\eps)N}\bigr)\to0$; hence
$\limsup_N\rho_k\le(k+1)H(\alpha)/\dF$.
\end{proposition}
\begin{proof}
Fix a support tuple $(S_0,\dots,S_k)$. The compressed product $M=F_k\cdots F_1$, with
$F_j=A_j[S_j,S_{j-1}]$ independent $n\times n$ Gaussian blocks, satisfies $\smin(M)\ge\prod_{j=1}^k\smin(F_j)$.
Writing $X_j=-\log\smin(F_j)\ge0$, the hard-edge tail \eqref{eq:hardedge} gives
$\PP(X_j\ge x)=\PP(\smin(F_j)\le e^{-x})\le C_\F(\sqrt n\,e^{-x})^{\dF}=C_\F n^{\dF/2}e^{-\dF x}$, so each
$X_j$ has an exponential upper tail of rate $\dF$ and prefactor $a=C_\F n^{\dF/2}$. The $X_j$ are
independent, so the sum $\Sigma=\sum_{j=1}^kX_j$ has, by the standard convolution bound for exponential
tails,
\[
  \PP(\Sigma\ge L)\ \le\ C_k\,a^k\,\frac{L^{k-1}}{(k-1)!}\,e^{-\dF L}\qquad(L\ge k\log a),
\]
by stochastic domination of each $X_j$ by a shifted exponential $\tfrac{\log a}{\dF}+\mathrm{Exp}(\dF)$
and the Gamma tail of their sum; this is the only range used, since $L=\Theta(N)$ below.
With $L=-\log t$ and $t=e^{-((k+1)H/\dF+\eps)N}$, this is $e^{o(N)}t^{\dF}$ (since $a^kL^{k-1}=e^{o(N)}$ for
fixed $k$, $n\le N$, $L=O(N)$), so $\PP(\smin(M)\le t)\le e^{o(N)}t^{\dF}$. A union bound over the
$\binom Nn^{k+1}=e^{(k+1)H(\alpha)N+o(N)}$ tuples gives
$\PP(M_N^{(k)}\le t)\le e^{(k+1)H(\alpha)N+o(N)}t^{\dF}=e^{-\dF\eps N+o(N)}\to0$.
\end{proof}

\subsection{A two-sided-singular factor: $\rho_k\ge 2H/\dF$}
\begin{proposition}\label{prop:prod-lower}
For every $\eps>0$, $\PP\bigl(M_N^{(k)}(\alpha)>e^{-(2H(\alpha)/\dF-\eps)N}\bigr)\to0$; hence
$\liminf_N\rho_k\ge2H(\alpha)/\dF$.
\end{proposition}
\begin{proof}
By \Cref{thm:twosided} applied to $A_1$, w.h.p.\ there are supports $S_0,S_1$ and a unit vector
$u\in\F^{S_0}$ (the least right singular vector of $F_1=A_1[S_1,S_0]$) with
$\|F_1u\|=\smin(F_1)\le e^{-(2H(\alpha)/\dF-\eps/2)N}$. Extend to any supports $S_2,\dots,S_k$. Then, with
$F_j=A_j[S_j,S_{j-1}]$,
\[
  M_N^{(k)}(\alpha)\ \le\ \smin(F_k\cdots F_1)\ \le\ \|F_k\cdots F_1u\|\ \le\
  \Bigl(\prod_{j=2}^k\|F_j\|_\op\Bigr)\|F_1u\|.
\]
Each $\|F_j\|_\op\le 2\sqrt n(1+o(1))$ w.h.p.\ by Bai--Yin (choosing, say, $S_j=[n]$; here the entries
are standard, so the edge is at scale $\sqrt n$), so the product of the $k-1$ operator norms is at most
$(2\sqrt n)^{k-1}(1+o(1))=e^{o(N)}$ for fixed $k$, and
$M_N^{(k)}\le e^{-(2H(\alpha)/\dF-\eps/2)N}e^{o(N)}\le e^{-(2H(\alpha)/\dF-\eps)N}$ w.h.p. The rate is
insensitive to the entry normalization, since scaling $A_j$ by a constant shifts $-N^{-1}\log\smin$ by
$o(1)$.
\end{proof}

This already improves the lower end of the bracket recorded in \cite{shmaloProducts} from $H(\alpha)/\dF$
(one one-sided-singular factor) to $2H(\alpha)/\dF$ (one two-sided-singular factor): the two-sided
construction of \Cref{app:twosided} is strictly better than making a single factor one-sided singular.

\subsection{The sharp constant and why the cascade is not enough}
\begin{remark}[Achievability of $(k+1)H/\dF$]\label{rem:achieve}
The numerics of \Cref{sec:numerics} determine $\rho_k=(k+1)H(\alpha)/\dF$: at $k=2$, $N=8$ the exhaustive
minimum over all support triples has rate $2.087$, on the predicted $3H(1/2)=2.079$, and the ratio
$\rho_2/\rho_1=1.53$ matches the value $1.5$ that $(k+1)H/\dF$ requires and excludes the $2H/\dF$ guess of
\cite{shmaloProducts} (ratio $1$). The following records what is and is not proved about the
achievability of this value.

The upper end $(k+1)H/\dF$ is a first-moment threshold: there are $e^{(k+1)H(\alpha)N}$ support tuples, and
a single fixed product $F_k\cdots F_1$, being a product of $k$ independent square Gaussian blocks, has the
product hard edge $\PP(\smin\le t)=t^{\dF+o(1)}$ (the extra factors contribute logarithms, not a power;
Akemann--Ipsen--Kieburg, Forrester), so the expected number of tuples with $\smin\le
e^{-((k+1)H/\dF-\eps)N}$ is $e^{(k+1)H(\alpha)N}e^{-((k+1)H(\alpha)-\dF\eps)N}=e^{\dF\eps N}\to\infty$. The
mean diverges below the threshold, so a matching second moment gives achievability.

The construction that realizes it is a cascade, but not the naive one. Suppose $S_0,\dots,S_{k-1}$ are
already chosen so that the partial product $Q=F_{k-1}\cdots F_1$ (an $n\times n$ matrix on rows $S_{k-1}$,
columns $S_0$) has $\smin(Q)\le e^{-(kH(\alpha)/\dF-\eps)N}$, with right singular vector $v\in\F^{S_0}$ (so
$\|Qv\|=\smin(Q)$). Adjoin the last factor by choosing the rows $S_k$: the composite
$M(S_k)=A_k[S_k,S_{k-1}]\,Q$ has rows $m_i=a_i^\top Q$ with $a_i=A_k[i,S_{k-1}]$ i.i.d.\ standard Gaussian,
so $M(\,\cdot\,)$ is the row-selection of an $N\times n$ matrix whose i.i.d.\ rows have covariance
$\Sigma=Q^\top Q$ (acting on $\F^{S_0}$). Minimizing $\smin(M(S_k))$ over row supports is therefore the
one-sided extremal least singular value problem of \cite{shmaloBAMS} for an \emph{anisotropic} Gaussian
row ensemble, and the marked-row construction there makes $\smin$ small at the one-sided rate
$H(\alpha)/\dF$ times the scale $\sqrt{\lambda_{\min}(\Sigma)}=\smin(Q)$ of the covariance's smallest
eigenvalue: $\Sigma$ inherits the tiny eigenvalue $\smin(Q)^2$ along $v$, and the row selection adds its
own $H(\alpha)/\dF$ of entropy on top. This yields $\smin(M(S_k))\le e^{-((k+1)H(\alpha)/\dF-\eps')N}$, one
more $H(\alpha)/\dF$ per layer,
telescoping from the two-sided base $\rho_1=2H(\alpha)/\dF$ to $\rho_k=(k+1)H(\alpha)/\dF$. The naive
cascade that fixes the carried direction $\eta$ and steers it --- minimizing $\|P_{S_k}A_k\eta\|$ over
$S_k$ --- gains only the $O(1)$ factor $\sqrt{\kcheck(\alpha)}$ and reaches merely $2H(\alpha)/\dF$; the
gain comes instead from re-optimizing the null direction of the whole composite, which the anisotropic
one-sided law does.

The first-moment side of this cascade step is rigorous, with no condition on $\Sigma$; only the
second moment remains. This is the content of \Cref{prop:aniso-mean} below, which shows that the
anisotropic marked count diverges in mean below the threshold $(k+1)H(\alpha)/\dF$. What obstructs a
full proof is the covariance $\Sigma=Q^\top Q$ being \emph{exponentially} ill-conditioned --- its
condition number is $\smax(Q)^2/\smin(Q)^2=e^{\Theta(N)}$ --- so the second-moment concentration cannot
be quoted from the subexponential-covariance regime of \cite{shmaloBAMS}, and must be re-run with the
configuration-dependent prefactor identified in \Cref{prop:aniso-mean}. The numerics of
\Cref{sec:numerics} confirm the value (the cascade reaches $(k+1)H/\dF$; at $k=2$, $N=8$ the greedy
cascade rate is $1.95$, approaching $3H=2.08$ from below while the exhaustive minimum sits on $3H$), and
the anisotropic scaling itself is checked directly in \texttt{verify\_anisotropic\_onesided.py}. The
bracket of \Cref{thm:product} is what is proved; the sharp value is the numerically determined answer
inside it.
\end{remark}

\begin{proposition}[First moment of the anisotropic cascade step]\label{prop:aniso-mean}
Let $M$ be $N\times n$ with i.i.d.\ rows $m_i\sim N(0,\Sigma)$, $\Sigma\succ0$ with least eigenvalue
$\lambda_{\min}$, and let $\mathcal M\subseteq[N]$ be a marked set with $|\mathcal M|=\lceil\sqrt N\rceil$.
For a configuration $c=(R,y)$ with $R\subseteq[N]\setminus\mathcal M$, $|R|=n-1$, and $y\in\mathcal M$, let
$u_R$ be a unit normal to $\{m_i:i\in R\}$ and $E_c=\{|\langle m_y,u_R\rangle|\le t\}$; set
$W=\sum_c\one_{E_c}$. Then:
\begin{enumerate}
\item[(i)] $W>0$ implies $\min_{|T|=n}\smin(M[T,\cdot])\le t$;
\item[(ii)] writing $\nu$ for the (Haar-distributed) unit normal to $n-1$ i.i.d.\ standard Gaussian
vectors, $\|\Sigma^{1/2}u_R\|=1/\|\Sigma^{-1/2}\nu\|$ in law, and $\PP(E_c\mid\{m_i\}_{i\in R})=p\bigl(t\,\|\Sigma^{-1/2}\nu\|\bigr)$
exactly, so
\[
  \EE W=\binom{N}{n-1}|\mathcal M|\;\EE_\nu\!\bigl[p\bigl(t\,\|\Sigma^{-1/2}\nu\|\bigr)\bigr];
\]
\item[(iii)] consequently $\EE W\ \ge\ c\,\binom{N}{n-1}|\mathcal M|\;\dfrac{t}{\sqrt{n\,\lambda_{\min}}}$
for a constant $c=c(\F)>0$ and all $t\le\sqrt{n\lambda_{\min}}$, so $\EE W\to\infty$ whenever
$t\ge\sqrt{n\lambda_{\min}}\,e^{-(H(\alpha)/\dF-\eps)N}$.
\end{enumerate}
Applied along the cascade with $\Sigma=Q^\top Q$ and $\lambda_{\min}=\smin(Q)^2=e^{-(2kH(\alpha)/\dF-o(1))N}$,
part (iii) gives $\EE W\to\infty$ for every $t\ge e^{-((k+1)H(\alpha)/\dF-\eps)N}$: the first moment
diverges below the sharp threshold.
\end{proposition}

\begin{proof}
(i) On $E_c$ the row set $T=R\cup\{y\}$ has $|T|=n$, and $\|M[T,\cdot]\,u_R\|^2=\sum_{i\in T}|\langle m_i,u_R\rangle|^2=|\langle m_y,u_R\rangle|^2\le t^2$
since $\langle m_i,u_R\rangle=0$ for $i\in R$; as $u_R$ is a unit vector, $\smin(M[T,\cdot])\le t$.

(ii) Write $m_i=\Sigma^{1/2}g_i$ with $g_i$ i.i.d.\ standard Gaussian. The condition $m_i\perp u_R$ for
$i\in R$ reads $\langle g_i,\Sigma^{1/2}u_R\rangle=0$, so $\Sigma^{1/2}u_R$ is parallel to the unit normal
$\nu$ of $\{g_i:i\in R\}$, which is Haar because the $g_i$ are isotropic. Thus $\Sigma^{1/2}u_R=c\nu$ with
$c=\|\Sigma^{1/2}u_R\|$, and $u_R=c\,\Sigma^{-1/2}\nu$; unit norm forces $c=1/\|\Sigma^{-1/2}\nu\|$. Since
$y\notin R$, $g_y$ is independent of $u_R$, and $\langle m_y,u_R\rangle=\langle g_y,\Sigma^{1/2}u_R\rangle=c\langle g_y,\nu\rangle$
is $N(0,c^2)$ given the defining rows, so $\PP(E_c\mid\{m_i\}_{i\in R})=p(t/c)=p(t\|\Sigma^{-1/2}\nu\|)$.
Taking expectations and multiplying by the number $\binom{N}{n-1}|\mathcal M|$ of configurations gives the
stated mean.

(iii) Diagonalize $\Sigma=\sum_j\lambda_j w_jw_j^\top$ with $\lambda_{\min}=\lambda_n$ the least eigenvalue
and eigenvector $w_n$. Then $\|\Sigma^{-1/2}\nu\|^2=\sum_j\lambda_j^{-1}\langle\nu,w_j\rangle^2\ge\lambda_n^{-1}\langle\nu,w_n\rangle^2$,
so $\|\Sigma^{-1/2}\nu\|\ge\lambda_n^{-1/2}|\langle\nu,w_n\rangle|$. Using $p(s)\ge c_\F\min(s,1)$ (valid for
every standard $\F$-Gaussian small-ball probability) and $t\le\sqrt{n\lambda_n}$,
\[
  \EE_\nu\bigl[p(t\|\Sigma^{-1/2}\nu\|)\bigr]\ \ge\ c_\F\,\EE_\nu\bigl[\min\bigl(t\lambda_n^{-1/2}|\langle\nu,w_n\rangle|,\,1\bigr)\bigr]\ \ge\ c_\F'\,t\,\lambda_n^{-1/2}\,\EE_\nu|\langle\nu,w_n\rangle|,
\]
where the last step keeps only the part of the range where the minimum equals its first argument, which
has $\Theta(1)$ probability. For Haar $\nu$ on $\Sph^{n-1}$, $\EE_\nu|\langle\nu,w_n\rangle|=\Theta(1/\sqrt n)$,
so $\EE_\nu[p(t\|\Sigma^{-1/2}\nu\|)]\ge c''\,t/\sqrt{n\lambda_n}$, and (iii) follows from (ii) with
$\binom{N}{n-1}|\mathcal M|=e^{H(\alpha)N+o(N)}$.
\end{proof}

\begin{remark}[The remaining second moment]\label{rem:aniso-second}
By \Cref{prop:aniso-mean}(i) and the second-moment method, achievability of the cascade step reduces to
$\Var(W)=o((\EE W)^2)$. Two configurations $c,c'$ with distinct marked rows are, unlike the isotropic
case of \Cref{app:twosided}, no longer exactly independent: the conditional probability
$p(t\|\Sigma^{-1/2}\nu\|)$ depends on the configuration through its normal $\nu$, so
$\PP(E_c\cap E_{c'})=\EE[\one_{E_{c'}}p(t\|\Sigma^{-1/2}\nu_c\|)]$ carries a covariance between the
prefactor $\|\Sigma^{-1/2}\nu_c\|$ and the event $E_{c'}$ whenever the defining row sets $R,R'$ overlap.
This covariance is the sole obstruction: it vanishes for disjoint $R,R'$ and is otherwise controlled by
the overlap of the two Haar normals, exactly the quantity the two-direction estimate
\Cref{lem:twoslab} bounds --- the anisotropy enters only through the scalar prefactor, since
$\Sigma^{1/2}u_R/\|\Sigma^{1/2}u_R\|=\nu$ is Haar regardless of $\Sigma$. Carrying the overlap-region
analysis of \Cref{app:twosided} through with this varying prefactor is the remaining step, and the one
this paper leaves open.
\end{remark}

\begin{remark}[What the network application needs]\label{rem:app-needs}
The certification-gap theorem (\Cref{thm:cert-gap}) uses only the \emph{proved} bracket: a pattern-uniform
certificate must tolerate the worst-conditioned pattern, whose least singular value is at most
$e^{-2H(\alpha)N/\dF}$ by the proved achievability of \Cref{prop:prod-lower} and at most
$e^{-(k+1)H(\alpha)N/\dF}$ at the sharp (numerically determined) rate. So the gap is proved to be at least
$e^{2H(\alpha)N/\dF}$ and equals $e^{(k+1)H(\alpha)N/\dF}$ at the sharp rate; the qualitative conclusion ---
that pattern-uniform certificates degrade exponentially in $N$ --- is unconditional.
\end{remark}

\section{Operator-norm brackets and the gate constants}\label{app:brackets}

This appendix proves the operator-norm brackets of \Cref{sec:rmt} and the order-statistic and
convex-hull facts used by the network estimates. Throughout, a Gaussian matrix has entry variance
$1/q$ where $q$ is the number of columns, unless a different normalization is stated.

\subsection{The order-statistic law}
\begin{lemma}[Order statistics of squared Gaussians]\label{lem:order}
Let $X_1,\dots,X_p$ be i.i.d.\ nonnegative integrable with continuous law and let $X^{\downarrow}_{1}\ge\cdots\ge X^{\downarrow}_p$
be the decreasing rearrangement. If $q/p\to\alpha\in(0,1]$ then, almost surely and in $L^1$,
\[
  \frac1p\sum_{j=1}^q X^{\downarrow}_j\ \to\ \EE[X\one\{X\ge\tau(\alpha)\}],\qquad
  \frac1p\sum_{j=p-q+1}^p X^{\downarrow}_j\ \to\ \EE[X\one\{X\le\tau'(\alpha)\}],
\]
where $\PP(X\ge\tau(\alpha))=\alpha$ and $\PP(X\le\tau'(\alpha))=\alpha$. Applied to $X_i=Z_i^2$ with
$Z_i\sim N(0,1)$ these limits are $\kappa(\alpha)$ and $\kcheck(\alpha)$ of \eqref{eq:kappa}--\eqref{eq:kcheck};
the deviation probability at any fixed threshold is $e^{-cp}$ by bounded-differences concentration.
\end{lemma}
\begin{proof}
Use the variational identity $\sum_{j\le q}X^{\downarrow}_j=\min_{\lambda\ge0}\bigl(q\lambda+\sum_i(X_i-\lambda)_+\bigr)$.
The empirical mean $p^{-1}\sum_i(X_i-\lambda)_+\to\EE(X-\lambda)_+$ uniformly in $\lambda$ on compacts by
the law of large numbers and monotonicity, and $q/p\to\alpha$; the minimizer is $\lambda=\tau(\alpha)$
(the $\alpha$-quantile), giving $\alpha\tau(\alpha)+\EE(X-\tau(\alpha))_+=\EE[X\one\{X\ge\tau(\alpha)\}]$.
The bottom sum is $\sum_j X^{\downarrow}_j-\sum_{j\le p-q}X^{\downarrow}_j$, the same law with $1-\alpha$.
Concentration: each summand changes the statistic by a bounded amount after truncation, so McDiarmid
applies with the exponential rate.
\end{proof}

\subsection{Operator-norm brackets}
\begin{proof}[Proof of \Cref{prop:A2}]
Write $M_1(\alpha)=\max_{|S|=|T|=\alpha N}\|P_TAP_S\|_\op=\sup\{\langle y,Ax\rangle:\|x\|=\|y\|=1,\ |\supp x|\le\alpha N,\ |\supp y|\le\alpha N\}$.
The two index sets are $\alpha N$-sparse unit spheres $\Sigma_{\mathrm{in}},\Sigma_{\mathrm{out}}$, whose
Gaussian widths are, by \Cref{lem:order}, $w(\Sigma_{\mathrm{in}})=\EE\sup_{x\in\Sigma_{\mathrm{in}}}\langle g,x\rangle=\sqrt{N\kappa(\alpha)}(1+o(1))$
and likewise $\sqrt{N\kappa(\alpha)}$ for the output. Chevet's inequality bounds
$\EE M_1(\alpha)\le\mathrm{rad}(\Sigma_{\mathrm{out}})\,w(\Sigma_{\mathrm{in}})/\sqrt N+\mathrm{rad}(\Sigma_{\mathrm{in}})\,w(\Sigma_{\mathrm{out}})/\sqrt N=2\sqrt{\kappa(\alpha)}$
with unit radii and entry variance $1/N$, and Gaussian concentration ($M_1$ is $1$-Lipschitz in the
Frobenius norm, variance $1/N$) removes the expectation, giving the upper bound $2\sqrt{\kappa(\alpha)}+o(1)$.
For the lower bounds, a fixed $\alpha N\times\alpha N$ submatrix already has $\|P_TAP_S\|_\op\to2\sqrt\alpha$
(Bai--Yin), and for a fixed unit $x_0$ on $S$, $Ax_0\sim N(0,I/N)$, whose top $\alpha N$ coordinates have
norm $\sqrt{\kappa(\alpha)}+o(1)$ by \Cref{lem:order}, so $\max_T\|P_TAx_0\|\ge\sqrt{\kappa(\alpha)}$.
\end{proof}

\begin{proof}[Proof of \Cref{prop:A4}]
Upper: submultiplicativity $\|M\|_\op\le\prod_j\|F_j\|_\op$ and then relax the shared supports to
independent per-factor maxima, each bounded by $2\sqrt{\kappa(\alpha)}+o(1)$ from \Cref{prop:A2}. Lower:
sequential planting. Fix a unit $x_0$ on $S_0$; at step $j$, $x_{j-1}$ is a fixed vector,
$A_jx_{j-1}\sim N(0,I/N)$, and choosing $S_j$ the top $\alpha N$ coordinates gains
$\sqrt{\kappa(\alpha)}$, with $x_j$ its normalization, a function of $A_1,\dots,A_j$ only, so the next
step is independent; this gives the greedy lower bound $\kappa(\alpha)^{k/2}\cdot(\ge\sqrt{\kappa(\alpha)})$.
Boosting the first step by a fixed submatrix top singular pair ($2\sqrt\alpha$ by Bai--Yin) and keeping
its left singular vector gives the $\max\{2\sqrt\alpha,\sqrt{\kappa(\alpha)}\}\kappa(\alpha)^{(k-1)/2}$
form; only one end can be boosted, since boosting fixes the adjacent output direction.
\end{proof}

\begin{proof}[Proof of \Cref{prop:contract}]
From \Cref{prop:A4}, $\max_{S_\bullet}\|M\|_\op\le(2\sqrt{\kappa(\alpha)})^k$, which tends to $0$ iff
$2\sqrt{\kappa(\alpha)}<1$, i.e.\ $\kappa(\alpha)<1/4$; the root of $\kappa(\alpha^*)=1/4$ is
$\alpha^*=0.042672$ numerically. The fixed-support threshold is the Fuss--Catalan edge
$\alpha^{k/2}\sqrt{(k+1)^{k+1}/k^k}<1$, i.e.\ $\alpha<(k^k/(k+1)^{k+1})^{1/k}$.
\end{proof}

\subsection{The subgaussian operator norm}
\begin{lemma}[Subgaussian bracket]\label{lem:subg}
Let $A$ be $N\times N$ with i.i.d.\ mean-zero, variance-$1/N$ entries of subgaussian norm $\|\sqrt N A_{ij}\|_{\psi_2}\le K$.
Then $\limsup_N M_1(\alpha)\le C_0K(\sqrt{\kappa(\alpha)}+\sqrt{\kappa(\alpha)})=2C_0K\sqrt{\kappa(\alpha)}$
w.h.p., with $C_0$ absolute. The lower bracket $\max\{2\sqrt\alpha,\sqrt{\kappa(\alpha)}\}$ is unchanged
and universal.
\end{lemma}
\begin{proof}
Write $M_1(\alpha)=\sup_{x\in\Sigma_{\mathrm{in}},y\in\Sigma_{\mathrm{out}}}\langle y,Ax\rangle$ and apply
the subgaussian bilinear comparison (a Chevet inequality for subgaussian matrices, e.g.\ via a
chaining/generic-chaining bound on the sparse spheres): $\EE\sup\langle y,Ax\rangle\le C_0K(\mathrm{rad}(\Sigma_{\mathrm{out}})w(\Sigma_{\mathrm{in}})+\mathrm{rad}(\Sigma_{\mathrm{in}})w(\Sigma_{\mathrm{out}}))/\sqrt N$,
with the same widths $\sqrt{N\kappa(\alpha)}$. Bousquet-type concentration for suprema of subgaussian
processes removes the expectation. The lower bound is a fixed submatrix and a fixed direction, both of
which only use the variance.
\end{proof}

\subsection{The convex-hull reduction}
\begin{lemma}[Gates in the hull]\label{lem:hull}
Let $W$ be a fixed matrix and $\mathcal D_\pi$ the set of diagonal matrices with entries in $[0,1]$ and
support pattern $\pi$. Then $\sup_{D\in\mathcal D_\pi}\|DWx\|=\max_{D\in\{0,1\}^\pi}\|DWx\|$ for every $x$,
and $\inf_{D\in\mathcal D_\pi,\,D\succeq D_0}\smin(DW)\ge\smin(D_0W)$ for $D_0\in\{0,1\}^\pi$. Consequently
every operator-norm upper bound and least-singular-value lower bound proved for $0/1$ gates holds for any
$[0,1]$-slope activation with the same support pattern.
\end{lemma}
\begin{proof}
$D\mapsto\|DWx\|^2=\sum_i D_{ii}^2(Wx)_i^2$ is convex in the diagonal entries, hence maximized at a vertex
$D\in\{0,1\}^\pi$; and $D\mapsto\smin(DW)$ is monotone on the order interval (increasing a nonnegative
diagonal entry cannot decrease any singular value, by the min--max characterization applied to $W^*D^2W$),
so the least singular value over $D\succeq D_0$ is at least $\smin(D_0W)$. The forward-pass and
gradient-lower-bound estimates use exactly these two extremal facts.
\end{proof}

\section{Network lemmas}\label{app:netlemmas}

This appendix completes the proofs of the uniform gain, gradient lower bound, and stability estimates of
\Cref{sec:gains}. The random-matrix inputs are \Cref{lem:order,lem:hull} and \Cref{prop:A2} of
\Cref{app:brackets}; the network structure is that of \Cref{sec:networks}.

\subsection{Uniform two-sided gain}
\begin{proof}[Proof of \Cref{thm:uniform-gain}]
Fix a unit $x\in V$. The vector $Wx$ has i.i.d.\ $N(0,1/q)$ coordinates, so $q\|Wx\|^2$ is a
chi-square with $p$ degrees of freedom and, by \Cref{lem:order} applied to $Z_i^2=(q^{1/2}(Wx)_i)^2$,
the top and bottom $\alpha p$ squared coordinates satisfy
\[
  \frac1q\sum_{i\in T_{\mathrm{top}}}(Wx)_i^2\to\tfrac pq\kappa(\alpha),\qquad
  \frac1q\sum_{i\in T_{\mathrm{bot}}}(Wx)_i^2\to\tfrac pq\kcheck(\alpha),
\]
with deviation probability $e^{-cp}$ at any fixed threshold. For an arbitrary $T$ of size $\alpha p$,
$\sum_{i\in T}(Wx)_i^2$ lies between the bottom and top sums, giving the two-sided bracket for that $x$
and that $T$, with the extremes attained by $T=T_{\mathrm{top}},T_{\mathrm{bot}}$. To make this uniform
over $x\in V$, take an $\eps$-net $\mathcal N$ of the unit sphere of $V$, $|\mathcal N|\le(3/\eps)^d$;
apply the bound at each net point with threshold slack $\eps$, and union bound: the failure probability
is $(3/\eps)^d e^{-cp}\to0$ since $d\log(1/\eps)=o(p)$. For a general $x$, pick $x'\in\mathcal N$ with
$\|x-x'\|\le\eps$; then $\|P_TWx\|$ and $\|P_TWx'\|$ differ by at most $\|W\|_\op\eps\le C\sqrt{p/q}\,\eps$,
and the top/bottom-$\alpha$ selections for $x$ and $x'$ differ on at most $O(\eps p)$ coordinates, whose
contribution is $O(\sqrt\eps)$; choosing $\eps=1/p$ makes both corrections $o(1)$. This gives the
uniform two-sided bracket and the sharpness.
\end{proof}

\subsection{Gradient lower bound}
\begin{proof}[Proof of \Cref{prop:gradlb}, details]
For a fixed input $x$, the symmetric-weight decoupling of \cite[Lemma 2.1]{bubeckCGT21} gives that
$\nabla F(x)^\top=w^\top D_L(x)W_L\cdots D_1(x)W_1$ has the same law as $w^\top\Sigma_LW_L\cdots\Sigma_1W_1$,
where $\Sigma_\ell=\diag(\xi_{\ell,1},\dots,\xi_{\ell,n_\ell})$ has entries $\xi_{\ell,i}$ i.i.d.\ with the
law of $\sigma_\ell'(Z)$ ($Z$ the preactivation), independent of the $W$'s. Set
$v_0=w$ and $v_\ell^\top=v_{\ell-1}^\top\Sigma_\ell W_\ell$; then
$\EE[\|v_\ell\|^2\mid v_{\ell-1}]=\EE[\Sigma_\ell]^2$-weighted $\|v_{\ell-1}\|^2$ times the layer variance,
and the pointwise nondegeneracy $\sigma_\ell'(r)^2+\sigma_\ell'(-r)^2\ge2c^2$ combined with the
signal-preserving variance gives $\EE[\|v_\ell\|^2\mid v_{\ell-1}]\ge c_0^2\|v_{\ell-1}\|^2$ for an
absolute $c_0>0$ (for ReLU, $\EE\xi_{\ell,i}^2=1/2$ and the variance-$2$ normalization make the per-layer
factor $\ge1-o(1)$). Iterating, $\EE\|\nabla F(x)\|^2\ge c_g^2\|w\|^2$ with $c_g=c_0^L$; concentration of
the quadratic form over the $\Theta(\prod_\ell n_\ell)$ Gaussian entries (Hanson--Wright) gives
$\|\nabla F(x)\|\ge\tfrac12c_g\|w\|$ with failure probability $e^{-c'\min_\ell n_\ell}$. Uniformity over
$x$ is by the $\eps$-net argument as in \Cref{thm:uniform-gain}, extended to all $x$ by
\Cref{thm:stability}. The nondegeneracy holds for any $[0,1]$-slope $\sigma$ with $\EE\sigma'(Z)^2>0$ by
\Cref{lem:hull} (the derivative gate lies in the hull of $0/1$ gates and the extremal least-singular
bound transfers); biases shift the gate density, keeping it in $[c_0,c_1]$ by \Cref{ass:nondeg}, which
only changes $c_0$.
\end{proof}

\subsection{Stability}
\begin{lemma}[Few gates flip]\label{lem:fewflip}
Under the hypotheses of \Cref{prop:gradlb}, w.h.p.\ for all inputs $x$ on the sphere and $x'$ with
$\|x'-x\|\le r$, the number of neurons at layer $\ell$ whose gate differs between $x$ and $x'$ is at most
$\delta_\ell n_\ell$ with $\delta_\ell=C r\sqrt{\bar n}/\|x\|$.
\end{lemma}
\begin{proof}
A neuron at layer $\ell$ flips only if its preactivation crosses a breakpoint of $\sigma_\ell'$, i.e.\ if
$z_{\ell,i}(x)$ lies within $\|z_\ell(x')-z_\ell(x)\|_\infty$ of a breakpoint. The layer maps are
$O(1)$-Lipschitz (upper bracket of \Cref{thm:uniform-gain}), so $\|z_\ell(x')-z_\ell(x)\|_\infty\le\|z_\ell(x')-z_\ell(x)\|_2\le Cr$
after the forward Lipschitz product. The preactivations $z_{\ell,i}$ are Gaussian of scale
$\varsigma_\ell\asymp\|h_{\ell-1}\|/\sqrt{n_{\ell-1}}\asymp\|x\|/\sqrt{n_{\ell-1}}$, so by anticoncentration
$\PP(|z_{\ell,i}-\text{breakpoint}|\le\tau)\le C\tau/\varsigma_\ell$ with $\tau=Cr$; the expected number of
flipped neurons is $n_\ell\cdot Cr\sqrt{n_{\ell-1}}/\|x\|\le C r\sqrt{\bar n}/\|x\|\cdot n_\ell$, and a
Bernstein bound gives the same w.h.p.\ uniformly over $x,x'$ by a net (the flip count is a Lipschitz
function of $x,x'$ after smoothing the indicator).
\end{proof}

\begin{proof}[Proof of \Cref{thm:stability}, details]
By the telescoping identity in the sketch, $\nabla F(x')^\top-\nabla F(x)^\top$ is a sum of $L$ terms, the
$\ell$-th being $w^\top D_L(x')W_L\cdots(D_\ell(x')-D_\ell(x))W_\ell\cdots D_1(x)W_1$. The gate difference
$D_\ell(x')-D_\ell(x)$ has entries in $[-1,1]$ supported on the flipped neurons, of number
$\le\delta_\ell n_\ell$ with $\delta_\ell=Cr\sqrt{\bar n}/\|x\|$ by \Cref{lem:fewflip}, so
$(D_\ell(x')-D_\ell(x))W_\ell$ restricted to that support has operator norm at most
$2\sqrt{\kappa(\delta_\ell)}+o(1)$ by \Cref{prop:A2} (a Gaussian block on a density-$\delta_\ell$ row
pattern), and $\kappa(\delta)\le C\delta\log(1/\delta)$ as $\delta\to0$ (the top-$\delta$ Gaussian
square-mass). The remaining factors are $O(1)$ in operator norm by \Cref{thm:uniform-gain}, and the
leading $w^\top D_L(\cdots)$ contributes $\|w\|$. Hence the $\ell$-th term is at most
$C\|w\|\sqrt{\delta_\ell\log(1/\delta_\ell)}$, and summing over $\ell$ (all $\delta_\ell\le\delta:=Cr\sqrt{\bar n}/\|x\|$)
gives $\|\nabla F(x')-\nabla F(x)\|\le C(L)\|w\|\sqrt{\delta\log(1/\delta)}$, which is the stated bound.
Uniformity over $x,x'$ is again by a net and the $O(1)$-Lipschitz dependence of every factor on the
inputs.
\end{proof}

\subsection{Combining into (A3)}
For the master theorem we need $\|\nabla F(x')-\nabla F(x)\|\le C_g/3$ on the attack ball of radius
$2a/C_g$. With $a=C\|x\|\sqrt{(d\log\bar n)/n_L}$ and $C_g=c_g\|w\|$, the radius is
$r=2a/C_g\asymp\|x\|\sqrt{(d\log\bar n)/n_L}/c_g$, so $\delta=Cr\sqrt{\bar n}/\|x\|\asymp\sqrt{(d\log\bar n)\bar n/n_L}/c_g$,
and $\|\nabla F(x')-\nabla F(x)\|\le C(L)\|w\|\sqrt{\delta\log(1/\delta)}=o(\|w\|)=o(C_g)$ whenever
$\bar n\,d\log\bar n=o(n_L\,c_g^2)$; when the widths are comparable the directional (attack-line) form of
\Cref{thm:stability} and the constant-$2$ Taylor engine of \cite[Thm 4.6]{advExt} require only
$\|\nabla F(x')-\nabla F(x)\|=o(\|\nabla F\|)$ along the segment, which holds under $d\log n=o(n)$
(\Cref{rem:widthcond}).


\begin{thebibliography}{99}

\bibitem{advExt}
Y.~Shmalo,
\emph{Constructive adversarial fragility of deep neural networks, ResNets, and vision transformers},
preprint, 2026.

\bibitem{bartlettBC21}
P.~L. Bartlett, S.~Bubeck, and Y.~Cherapanamjeri,
\emph{Adversarial examples in multi-layer random ReLU networks},
Advances in Neural Information Processing Systems 34 (2021).

\bibitem{bubeckCGT21}
S.~Bubeck, Y.~Cherapanamjeri, G.~Gidel, and R.~Tachet des Combes,
\emph{A single gradient step finds adversarial examples on random two-layer neural networks},
Advances in Neural Information Processing Systems 34 (2021).

\bibitem{bubeckSS19}
S.~Bubeck, Y.~T. Lee, E.~Price, and I.~Razenshteyn,
\emph{Adversarial examples from computational constraints},
Proceedings of the 36th International Conference on Machine Learning (2019), 831--840.

\bibitem{danielyM26}
A.~Daniely and E.~Mehalel,
\emph{Most convolutional networks suffer from small adversarial perturbations},
preprint arXiv:2602.03415, 2026.

\bibitem{danielyS20}
A.~Daniely and H.~Shacham,
\emph{Most ReLU networks suffer from $\ell^2$ adversarial perturbations},
Advances in Neural Information Processing Systems 33 (2020), 6629--6636.

\bibitem{edelman}
A.~Edelman,
\emph{Eigenvalues and condition numbers of random matrices},
SIAM J. Matrix Anal. Appl. 9 (1988), 543--560.

\bibitem{fawzi18}
A.~Fawzi, H.~Fawzi, and O.~Fawzi,
\emph{Adversarial vulnerability for any classifier},
Advances in Neural Information Processing Systems 31 (2018).

\bibitem{goodfellow15}
I.~J. Goodfellow, J.~Shlens, and C.~Szegedy,
\emph{Explaining and harnessing adversarial examples},
International Conference on Learning Representations (2015).

\bibitem{relumap}
Y.~Shmalo,
\emph{Finite-horizon decay of iterated ReLU maps with a fixed Gaussian matrix},
preprint, 2026.

\bibitem{rv}
M.~Rudelson and R.~Vershynin,
\emph{The Littlewood--Offord problem and invertibility of random matrices},
Adv. Math. 218 (2008), 600--633.

\bibitem{shafahi19}
A.~Shafahi, W.~R. Huang, C.~Studer, S.~Feizi, and T.~Goldstein,
\emph{Are adversarial examples inevitable?},
International Conference on Learning Representations (2019).

\bibitem{shmaloAccuracy}
L.~Berlyand, E.~Sandier, Y.~Shmalo, and L.~Zhang,
\emph{Enhancing accuracy in deep learning using random matrix theory},
preprint arXiv:2310.03165, 2024.

\bibitem{shmaloAdv}
Y.~Shmalo,
\emph{Adversarial properties of random deep neural networks via balance and spherical concentration},
preprint, 2026.

\bibitem{shmaloBAMS}
Y.~Shmalo,
\emph{Extreme least singular values with redundancy, radial ensembles, and quantitative stability of phase retrieval},
Bull. Amer. Math. Soc., to appear (2026).

\bibitem{shmaloProducts}
Y.~Shmalo,
\emph{Extreme singular values of products of Gaussian submatrices},
preprint, 2026.

\bibitem{shmaloRMT}
Y.~Shmalo, J.~Jenkins, and O.~Krupchytskyi,
\emph{Deep learning weight pruning with RMT-SVD: increasing accuracy and reducing overfitting},
preprint arXiv:2303.08986, 2023.

\bibitem{shmaloSharp}
Y.~Shmalo,
\emph{The sharp exponent for the smallest singular value of a coordinate-optimized Gaussian submatrix},
preprint, 2026.

\bibitem{szegedy14}
C.~Szegedy, W.~Zaremba, I.~Sutskever, J.~Bruna, D.~Erhan, I.~Goodfellow, and R.~Fergus,
\emph{Intriguing properties of neural networks},
International Conference on Learning Representations (2014).

\end{thebibliography}


\end{document}
